\chapter{Optimization, or Going to Extremes}
% \baselineskip=1.5cm
\label{cha:what-deriv-tells}
\section*{Introduction}
\label{sec:introduction}
\markboth{\sc Optimization}{Introduction}

While it would not be entirely correct to say that Calculus was
invented to solve optimization problems, it would not be entirely
wrong either.  Recall that the opening phrase of the title of
Leibniz's $1684$ paper on Differential Calculus was \pubtitle{``A
  New Method for Maxima and Minima .\ .\ .}.'' Also recall that a
fundamental tenet of modern science is that {\sc{}Nature Is
  Lazy}. That is, nature optimizes. Light follows the \alert{fastest}
(minimal) path between two points. Soap bubbles are spherical because
that is the shape that encloses the \alert{most} (maximal) volume with
the a given surface area. Alternatively, it uses the \alert{least}
(minimal) surface area to enclose a particular volume.

Mathematically it all comes to the same point: We have some quantity
(revenue, tax liability, temperature, resource consumption, travel
time, potential energy, volume, surface area, whatever), which is
(usually) constrained in some way, and we need to find the conditions
that make it as large, or as small as possible within those
constraints.

Usually the ``thing to be optimized'' is called the \term{Objective
  Function} and the constraints on the variables are called, well,
they're called the \term{Constraints}.

In chapter~\ref{cha:newt-metaph-motion} we observed that the optimal\aside{To
  ``optimize'' is to find the optimum value, whether that be a maximum
  or a minimum.} points on a given curve will occur where the curve is
horizontal (equivalently, where the derivative is zero). So it
seems that a likely strategy for attacking these kinds of problems
would be to express the objective function, $f(x)$,  in terms of the  variables
and then solving the equation $\dfdx{f}{x} = 0$. At their most
fundamental level optimization problems really are just that
straightforward.

But ``straightforward'' does not mean easy.

Recall that when we examined  Fermat's  ``Method of Aedequality'' in
section~\ref{sec:ferm-meth-aedeq} we said that, ``.\ .\ . it is a general
property of the graphs of functions that near a maximum or minimum:
$f(x+h)\approx f(x)$ as long as $h$ is not too big.'' A moment's
thought should convince you that this really says that at a maximum or
a minimum the graph of a function is practically horizontal. That is
why the following is called Fermat's Theorem:

\begin{mytheorem}[Fermat's Theorem]
  \label{thm:FermatsTheorem}
  If a function, $f(x)$ is differentiable at $x=a$ and
  $f(a)$ is either a maximum or a minimum then
  $\eval{\dfdx{f}{x}}{x}{a}=0.$
\end{mytheorem}

Take careful notice of what Fermat's Theorem says and, just as
important, what it does \emph{not} say: It says that if a point on
the graph of $y=f(x)$, say $(a, b)$ is \emph{known} to be either a maximum or a
minimum it must be that $\eval{\dfdx{y}{x}}{x}{a}=0$. It does
\alert{not} say that if $\eval{\dfdx{y}{x}}{x}{a}=0$ then $f(a)$ is
either a maximum or a minimum. In fact, that is not necessarily true.

Fermat's Theorem is given in terms of the derivative, $\dfdx{f}{x}$,
rather than differentials, as modern Calculus statements should.
However, if we think in terms of differentials it is clear that
setting $\dfdx{f}{x}=0$ is equivalent to setting $\dx{f}=0$, since a
fraction is zero precisely when its numerator is zero. We will use
whichever version is appropriate to the situation.  Either way, we
have a systematic way to approach optimization problems compared to
ad-hoc methods developed before of Newton and Leibniz.

\section{Some Simple Optimizations: Optimal Squares}\label{sec:five-simil-probl}
\markboth{\sc Optimization}{\sc Optimization Examples}
\label{example:LineDist}

In this section we will work carefully through several very similar
examples. All of these problems can be solved by methods other than
the optimization methods we will be using. We are just starting with
elementary problems to get a feel for optimization.


\begin{myexample}[Constructing A Square on a Line]
  \label{example:SqOnLine}

  Find the area of the smallest square that can be constructed with
  one corner at the point $(2,3)$ and one corner on the line
  $5x+3y=2$.

  Before we start, note that this problem is easy to visualize and,
  once visualized, it is intuitively clear that there is a square of
  minimim area. That is, the problem has a solution. It is worth
  noticing this explicitly. Not all of them
  do.

  In any optimization problem the first question we need to ask is,
  ``What do we need to optimize?'' In other words, ``What is our
  objective function?''

  \begin{wrapfigure}[]{r}{1.7in}
    \captionsetup{labelformat=empty}
    \centerline{\includegraphics*[height=1.6in,width=1.7in]{../Figures/LineDist1}}
    \label{fig:LineDist1}
  \end{wrapfigure}
  For this problem we want to find the area of the smallest square
  constructed in the manner shown in the sketch at the right. The
  length of the line segment between the point $(2,3)$ and \emph{any
    point on the line} $(x,y)$ is given by:$\sqrt{(x-2)^2+(y-3)^2}$.
  Therefore the area of the square constructed from this line segment
  is given by the function:
  $$
  A(x)=(x-2)^2+(y-3)^2,
  $$
  This is our objective function.

  Notice that we've abused our notation a bit. The notation $A(x)$
  clearly indicates that $A$ depends on $x$ alone but we've given $A$
  in terms of both $x$ and $y$. It should be clear that $A$ really
  depends only on $x$, since we can substitute $\frac13(2-5x)$ for $y$
  at any point.  Using two variables makes this objective function a
  bit ``easier on the eyes.''
  
  We need to find the smallest possible value of $A$, but the only
  tool we have so far is our observation in
  Section~\ref{sec:newtons-view:-motion} that a function has a
  horizontal tangent line -- and therefore attains an optimum value
  (either a maximum or a minimum) when its derivative is zero.

  Differentiating both sides of
  \begin{align*}
    A&=(x-2)^2+(y-3)^2\\
    \intertext{gives}
    \dx{A}&=2(x-2)\dx{x} + 2(y-3)\dx{y}.
            \intertext{Setting $\dx{A}=0$ we have}
            \dx{A}=2(x-2)\dx{x} + 2(y-3)\dx{y}=0.
  \end{align*} 
  But recall that $x$ and $y$ are
  \term{constrained} by the equation $5x+3y=2$. If we
  differentiate this we have $\dx{y}= -\frac{5}{3}\dx{x}$.
  Therefore
  \begin{align*}
    \dx{A}&= 2(x-2)\dx{x} + 2(y-3)\left(-\frac{5}{3}\dx{x}\right)\\
    \dx{A}&= \left[2(x-2) +
            2(y-3)\left(-\frac{5}{3}\right)\right]\dx{x}\\
    \intertext{or}
    \dfdx{A}{x} &= 2(x-2) +
                  2(y-3)\left(-\frac{5}{3}\right).\\
    \intertext{Setting $\dfdx{A}{x}$ equal to zero we have}\\
    2(x-2) + 2(y-3)\left(-\frac{5}{3}\right) &= 0.\\
    (x-2) + (y-3)\left(-\frac{5}{3}\right) &= 0.\\
    x-2-\frac53y+5&=0.\\
    x+3-\frac53y&=0.\\
    \intertext{From our \alert{constraint}, $5x+3y=2$, we have $y=\frac13(2-5x)$. Substituting gives}
    x+3-\frac53\left(\frac13(2-5x)\right)&=0,\\
    x+3-\frac53\left(\frac23-\frac{5}{3}x)\right)&=0,\\
    x+3-\frac{10}{9}+\frac{25}{9}x&=0,\\
    \frac{9x+27-10+25x}{9}&=0,\\
    \frac{34x+17}{9}&=0,\\
    \intertext{and solving for $x$ gives}
    34x&=-17\\
    x&=-\frac12.
  \end{align*}
  
  To find $y$ we again use our \alert{constraint}:
  \begin{align*}
    5x+3y&=2\\
    3y&=2-5x\\
    y&=\frac13(2-5x)\\
    \bigeval{y}{x}{-1/2} &= \eval{\frac13(2-5x)}{x}{-1/2}\\
    \bigeval{y}{x}{-1/2} &=\frac32.
  \end{align*}

  So we have found that the minimal square occurs when $x=-\frac12$,
  and $y=\frac32$. We're done, right?

  No, of course not. The problem asked us to find the \emph{area} of
  the minimal square. We've only found the \emph{location} of the
  minimal square.

  To find the area we need to put these numbers into our area formula:
  \begin{align*}
    \bigeval{A(x)}{\left(x,y\right)}{\left(-\frac12,\frac32\right)}&=\bigeval{(x-2)^2+(y-3)^2}{(x,y)}{\left(\frac{1}{2},\frac32\right)}\\
        &=\left(-\frac12-2\right)^2+\left(\frac12-3\right)^2=\frac{17}{2}\text{ square units.}
  \end{align*}

  This highlights an important aspect of optimization problems.  To
  solve them you need to first abstract a mathematical model from the
  physical problem by finding an appropriate objective function. But
  solving the model problem is not the same as solving the original
  problem. Once the model problem is solved you need to translate that
  into the context from which it came. In this case that means finding
  the area, not just $x$ and $y$.

A few more general comments are in order:

  \centerline{\bf{}\underline{Comments on  Example~\ref{example:LineDist}}}
  \begin{enumerate}
  \item Notice that without an objective function we can't even get
    started. The first step in any optimization problem is \emph{always}
    to identify the objective function.
    \begin{embeddedproblem}
      The objective function used in example~\ref{example:SqOnLine}
      assumed that $(2,3)$ and $(x,y)$ were adjacent corners of the
      square.  This was never stated in the problem. What would the
      objective function look like if the points were diagonally
      opposite each other.  How would this change the problem?
    \end{embeddedproblem}
  \item We need to identify our constraint(s) as well, if they
    exist. (Not all optimization problems will have constraints.)  We
    don't need these immediately, but we will need them sooner or later.
  \item It is intuitively clear that there actually \emph{is} a
    minimal square; that a solution exists. All you have to do is draw
    a picture to see it.  This is so clear that you can probably get
    close to the solution by simply sketching the problem. This will
    be true often enough that it will be easy to start believing that
    you can rely on your intuition alone to solve optimization
    problems.

    \centerline{\bf{} This is a mistake.}

    Don't misunderstand us. Intuition is important. Very
    important. It is especially so when you first start thinking about
    a new problem. Use your intuition.

    But intuition is based on past experience. It is a way of saying
    that you expect today to go pretty much the same way yesterday
    did. Most of the time this is true, but you can't always count on
    it. Sometimes things don't go the way you expect and you never
    know when things will go awry so you \alert{must} verify that your
    intuition by some other means. (Calculus, in this case.)

    \begin{embeddedproblem}
      For example, it seems fairly intuitively clear that the minimal
      square occurs when the line through $(2,3)$ is orthogonal
      (perpendicular) to the graph of the line $5x+3y=2$.
      \begin{enumerate}[label={\bf{}(\alph*)}]
        \item Verify that our intuition is correct.
        \item Is it also true if the point on the line is not adjacent
          to $(2,3)$?
      \end{enumerate}
    \end{embeddedproblem}
    
  \item  There is a strong tendency to think of the algebraic
    manipulations in an optimization problem as ``just algebra'' and to
    be careless with them.

    \centerline{\alert{This is also a mistake.}}

    As you can see, most of the heavy lifting in this problem was in
    the Algebra, not the Calculus. This will be true of many (not all)
    of the mathematics problems you encounter from now on, and
    \emph{every step} in a given problem is just as important as every
    other step, whether you call it Algebra, Calculus, or Arithmetic.
    Most of the mathematics problems you encounter from now on will
    require skillfull use of at least two or three of these essential
    tools.


  \item     We were very careful to write out every single algebraic and
    arithmetic step in this problem in detail. We will not do this in
    the future. Frequently we will do several computations at once. We
    will do this for two reasons:
    \begin{enumerate}
    \item You can see how much space it took for us to fill in all of
      the details in this problem. If we do this for every elementary
      computation this text would be too long to be useful. We don't
      want that to happen.
    \item You need to be able to follow these kinds of computations and
      you will not learn to do so unless we give you the chance to
      practice the skill. So we will regularly do several elementary
      algebraic steps at once.
      
      But ``elementary'' does not mean ``easy.'' You will not always be
      able to follow them in your head, especially at first. So keep a
      writing tablet handy at all times. If you do not see how to get
      from one line of computation it is very natural to  just skip
      it.

\centerline{      \alert{Don't.} \alert{Do.} \alert{That.}}

Take out that pad and do the computation yourself.  If you do this you
will find that following the computations gets easier over time. If
you do not it will remain difficult
forever.\\
\centerline{\alert{We're serious.}}  \centerline{\alert{Do all
    elementary computations yourself!}}  It is
\alert{\underline{your}} responsibility to make sure you understand
the computations we show you.
    \end{enumerate}
  \end{enumerate}

\end{myexample}

\begin{myexample}[Constructing A Square on a Circle]
  \label{example:CircleDist}
  Find the area of the largest square that can be constructed with
  one corner at the point $(2,3)$ and one corner on the unit circle:
  $$
  x^2+y^2=1.
  $$
  We have the same objective function as
  before: $A(x)=(x-2)^2+(y-3)^2$. 
  So, also as before,
  \begin{equation}
  \dx{A} = 2(x-2)\dx{x} + 2(y-3)\dx{y}.\label{eq:ConstObj1}
  \end{equation}

  
  But this time our constraint is $x^2+y^2=1$ so
  $\dx{y}=-\frac{x}{y}\dx{x}$.  Plugging this into
  equation~\ref{eq:ConstObj1} we have
  \begin{align*}
    \dx{A} &= 2(x-2)\dx{x} + 2(y-3)\left(-\frac{x}{y}\right)\dx{x}\\
    \dfdx{A}{x}&= 2\left(-2+\frac{3x}{y}\right)\\
  \end{align*}
  Setting $\dfdx{A}{x}=0$ and solving we see that
  \begin{align}
    \nonumber    3x&=2y.\\
    \intertext{and since $y=\pm\sqrt{1-x^2}$ we have}
    \label{eq:Sqrt}3x&=\pm2\sqrt{1-x^2}.\\
    \intertext{The square root symbol, and the $\pm$, will make solving
    equation~\ref{eq:Sqrt} problematic, so
    let's get rid of them by squaring both sides of
    equation~\ref{eq:Sqrt}. this gives}
    \nonumber 9x^2&=4(1-x^2)\\
    \intertext{or}
    \label{eq:LostSolution}   13x^2&=4.
  \end{align}
  
  \begin{wrapfigure}[]{l}{2.5in}
    \captionsetup{labelformat=empty}
    \centerline{\includegraphics*[height=2in,width=2.5in]{../Figures/CircleDist3}}
    \label{fig:CircleDist3}
  \end{wrapfigure}

  So $x=\frac{2}{\sqrt{13}}$ and
  $y=\sqrt{1-\frac{4}{13}}=\frac{3}{\sqrt{13}}$.  Thus the area of
  our square is
  is
  $$A=\left(\frac{2}{\sqrt{13}}-2\right)^2+ \left(\frac{3}{\sqrt{13}}-3\right)^2
  \approx 6.789.$$


  But wait a second! Our intuition is complaining. Something is not right here. Do you see the
  difficulty?

  A sketch of this problem is given at the left.  The black square is
  the one that solves the problem and it is the one we found, but it
  clearly isn't the \emph{largest} square. The red square is clearly
  larger.

  We made a mistake when we solved equation~\ref{eq:LostSolution}. See
  if you can find it before reading on.

  \begin{Digression}[Absolute Value]

    Let's fill in the missing steps to see what happened. We said that
    the solution of $13x^2= 4$ is $x=\frac{2}{\sqrt{13}}$ but that is
    not true.  Do you see the problem?  You've encountered it before
    but in our experience most students are still puzzled by this. It
    is worth digressing a bit to clear this up once and for all.

    When asked to think of a ``number'' most people tend to think of
    \emph{positive} numbers. But there are just as many negative
    numbers as positive. Our preference for positive numbers is a bias
    borne of early childhood experience with arithmetic, where only
    the positive numbers are used.  Going forward you will need to
    overcome this bias by making a deliberate effort not to assume
    that the expressions you encounter are positive. You must
    consciously, and deliberately consider the possibility of negative
    numbers.

    For example, the steps to solving equation~\ref{eq:LostSolution}
    are as follows:
    \begin{align}
      \nonumber    13x^2&=4\\
      \nonumber    x^2&=\frac{4}{13}\\
      \label{eq:LostSolution2}    \sqrt{x^2}&=\frac{2}{\sqrt{13}}
    \end{align}

    If, in the next step, we write $x=\frac{2}{\sqrt{13}}$ then
    we have
    implicitly assumed that $x$ is a positive number.  Here's why.

    The square root symbol, $\sqrt{\ }$, \alert{always} means the
    \alert{positive} square root so both sides of
    equation~\ref{eq:LostSolution2} are positive. So far, so good. But
    if $\sqrt{x^2}$ is positive does that mean that $x$ has to be
    positive?

    Of course not.

    It \emph{can} be positive, but it doesn't \emph{have} to be
    positive. For example, if $x=2$ then $\sqrt{2^2}=\sqrt{4}$ and
    since the radical \emph{always} means the positive root we
    have $$\sqrt{2^2}=\sqrt{4}=2.$$

    But now suppose $x=-2$. In this case we get
    $ \sqrt{(-2)^2}=\sqrt{4}=2$ again, since the radical always gives
    a positive value. Notice that there is nothing special about
    $2$. If $x$ is any positive number then $\sqrt{x^2} = x$ but if
    $x$ is any negative number then $\sqrt{x^2}=-x$.

    \alert{Huh?!} How did we get a negative?

    \alert{We didn't!} Because we all have a strong bias toward
    positive numbers we automatically assumed that $x$ is positive,
    making $-x$ negative. But remember that $x$ was a negative number so $-x$
    must be a positive number. 
    
    We can tie all of this up notationally as follows:
    \begin{equation}
      \sqrt{x^2}=
      \begin{cases}
        x& \text{if } x\ge 0\\
        -x& \text{if } x< 0
      \end{cases}\label{eq:SqrtSqr}  
    \end{equation}
    So it appears that the function $f(x)=\sqrt{x^2}$ will always return
    the ``positive version'' of $x$, right?

    Wait a moment. Don't we already have a name
    for the function that does that?

    Sure we do. We call that the Absolute Value function.

    So it appears that $\sqrt{x^2}=\abs{x}.$ Remember this and you will
    never make the mistake we made in this example. Solving $13x^2=4$  we get
    \begin{align*}
      \sqrt{x^2}&=\frac{2}{\sqrt{13}}\\
      \abs{x}&=\frac{2}{\sqrt{13}}
    \end{align*}
    and the two numbers whose absolute values are equal to
    $\frac{2}{\sqrt{13}}$ are $x=\frac{2}{\sqrt{13}}$ and
    $x=-\frac{2}{\sqrt{13}}.$
  \end{Digression}

  % We even saw this when  we solved our constraint, $x^2+y^2=1$
  % to get $y=\pm\sqrt{1-x^2},$ but when we squared both sides of
  % equation~\ref{eq:LostSolution2} that problem went away.

  \begin{ProficiencyDrill}
    \label{problem:CircleDist2}
    Fill in the details of the solution of
    equation~\ref{eq:LostSolution} to show that there are two
    solutions to the problem in this example. You can easily see from
    the digram above which is minimal, and which is maximal, but can
    you tell \emph{from your solution procedure} which one is which?
  \end{ProficiencyDrill}
  
  \centerline{\underline{\bf{}Comments on Example~\ref{example:CircleDist}}}
  \begin{enumerate}
  \item Remember that you will usually have some physical, or
    geometric intuition about the optimization problems that you'll
    encounter. Use it. After all, It was our intuition that led us to
    the difficulty with, and better understanding of, the expression
    $\sqrt{x^2}$.But you want to be able to verify your intuition. If
    you can, make a guess based on your intuition before you start.
    Even if your intuition isn't very strong, you'll usually be able
    to tell whether you are expecting a positive or a negative
    solution, or whether the solution will be large or small. Think
    about what you expect to find \emph{before} you start working on
    the problem.

    When you've finished, compare your intuitive answer with the one you
    actually get. If they don't align then something is
    wrong. It might be that you've made a mistake, but it also might be
    that your intuition was wrong. Find out which it is. \alert{Don't}
    just give up and move on. Resolve the discrepancy. This will
    improve either your intuition or your problem solving skills.
  \item Notice how easy it was to miss the maximal square before we
    drew our sketch.  Drawing a sketch is an aid to your intuition and
    is a \emph{very} good habit. Cultivate it. Make this part of how
    you begin any problem.

  \end{enumerate}
\end{myexample}

\begin{ProficiencyDrill}
  \label{problem:CircleDist3}
  Show that the intersection of the line joining the point $(2,3)$ to
  the center of the circle and the circle itself are the two optimal
  points obtained before.  How does this relate to the previous
  problem~\ref{problem:CircleDist2}? Do you see, intuitively, that
  once again the line joining the point to the curve is orthogonal
  (perpendicular) to  (line tangent to) the curve?
\end{ProficiencyDrill}

Let's return to finding extremal squares on curves. So far we've placed
one corner at the point $(2,3)$. Suppose we try to
construct an optimal square on an ellipse this time. Since a circle is
a special form of an ellips this turns out to be, essentially, the same
problem as Problem~\ref{problem:CircleDist2}, with trickier algebra.

\begin{embeddedproblem}[Extremal Squares on an Ellipse]
  \label{problem:EllipseDist}
  There are two
  extremal squares with one corner at $(2,3)$ an the other on the ellipse:
  \begin{equation}
    x^2+\frac14y^2=2.\label{eq:EllipseConstraint}
  \end{equation}
  The problem is to find the area of
  both of them.
  \begin{enumerate}[label={\bf{}(\alph*)}]
  \item Draw a sketch of this problem.\\\comment{Human beings are visual creatures. We move through the
      world relying primarily on our sense of sight. It is thus through
      vision that we have our strongest intuitive understanding of the
      world. Allowing yourself to ``see'' a problem is a big step toward
      intuitive understanding. You should \emph{always} sketch a problem out
      as best you can to allow yourself to ``see'' it. This is an excellent way
      to engage your intuition.}
  \item 
      The objective function for this problem is again:
$$
      A(x) = (x-2)^2+(y-3)^2.
$$

    What is the constraint?
  \item Use the constraint to show
    that $$\dx{y} = -\frac{4x}{y}\dx{x}$$
    and that
    $$
    y=\pm2\sqrt{2-x^2}
    $$
    Apparently we get to choose whether we want the positive or the
    negative value of $y$. Does this make sense to you?
  \item Use the results in parts (b), and (c) to show that if we
    choose $y=2\sqrt{2-x^2}$ then
    $$\dfdx{A}{x}=-6x-4+\frac{12x}{\sqrt{2-x^2}}.$$
    What is $\dfdx{A}{x}$ if we choose $y=-2\sqrt{2-x^2}.$
  \item Use Newton's Method to show that if we choose $
    y=2\sqrt{2-x^2}$ then the coordinates of the only
    root of $\dfdx{A}{x}=0$ are $x\approx -0.273$ and
    and $y\approx -2.7755$.
  \item Use Newton's Method to show that if we choose $
    y=-2\sqrt{2-x^2}$ the the coordinates of the only
    root of $\dfdx{A}{x}=0$ are $x\approx 0.857$ and
    and $y\approx 2.25$.
  \item Approximate the  area of the  square associated with both of
    these points? Which one is the maximal square? 
  \end{enumerate}
\end{embeddedproblem}
Do you see that, just like in Problem~\ref{problem:CircleDist2} we
will always get both a maximal and a minimal square as in the
following sketch:\\
\centerline{\centerline{\includegraphics*[height=2in,width=3.1in]{../Figures/EllipseWithSquares}}}

Notice that for both squares one of the two sides that meet at the corner that actually lies on the
 ellipse appears to be tangent to the ellipse. Do you suppose it
 actually is tangent, or is this just an artifact of this particular
 problem?

 \begin{ProficiencyDrill}
   Find an equation of the line tangent to the ellipse at each of the
   points you found in Problem~\ref{problem:EllipseDist} and show that
   the apparent tangency is not an artifact. In fact those sides
   actually are tangent.
 \end{ProficiencyDrill}

 Since we've constructed squares, it follows immediately that the
 lines from $(2,3)$ to the ellipse are orthogonal (perpendicular) to
 the ellipse. A look back shows that this was also true of the square
 constructed on a line (Example~\ref{example:SqOnLine}) and a square
 constructed on a circle (Example~\ref{example:CircleDist}).

 Do you suppose this is a general phenomenon? Watch for it in the next
example.

\begin{embeddedproblem}
  From problem~\ref{problem:CircleDist2} to problem~\ref{problem:EllipseDist} we
  changed our constraints. For this problem we will change the
  objective function.

  Consider the following picture of the unit circle $x^2+y^2=1$ and
  the point\notetoself{Perhaps we should change this point to $(2,3)$
    just to not change too much at once?} $(-4,-3)$.\\
  \centerline{\includegraphics*[height=1.5in,width=2in]{../Figures/Lagrange1}}
  Notice that if we draw a line connecting $(-4,-3)$ to an arbitrary
  point $(x,y)$, on the circle, then there is a point where the slope
  of the line is maximum and a point where the slope is minimum. These
  are the outermost lines depicted in the figure.

  We wish to find an equation for each of these lines.
  \begin{enumerate}[label={\bf{}(\alph*)}]
   \item Find an appropriate  objective function for this problem.
   \item What is the constraint for this problem?
   \item Differentiate both your objective function and your
     constraint to show that the points where the maximum and minimum
     slope occur must satisfy
     \begin{equation}
     x^2+y^2=-4x-3y.\label{eq:MaxMinSlope}
        \end{equation}
     \comment{It is possible that you didn't come up with the same
       objective function we did. In that case you may have a
       different equation to solve. That's ok. Don't let it stop
       you. There is usually more than one way to solve a given
       problem. Different people will find different methods. Find the
       one that makes sense to you. }
   \item Solve equation~\ref{eq:MaxMinSlope} simultaneously with the
     constraint to find the equations of the lines with maximal and
     minimal slope.  To verify your answer graph both lines and the
     unit circle.
   \end{enumerate}
 \end{embeddedproblem}

 % There is more to learn from finding extremal squares. We will return
 % to it, but first let's look at the general technique of optimizing an
 % objective function with a constraint. To illustrate we'll look at a
 % familiar example.
 \begin{myexample}[Derivatives vs. Differentials]
   In chapter~\ref{chapter:ScienceBeforeCalc} we used Fermat's Method of
   Aedequality to show that out of all rectangles with a fixed
   perimeter, the square encloses the largest area. To use Calculus we
   need to identify an objective function and the  constraints for
   this problem.

   The objective function is easy. The area, $A$, of any rectangle is
   its length, $l$, times its width, $w$, so
   \begin{equation}
     L=lw\label{eq:SquareObjFun}
  \end{equation}
  is the objective function.

  The perimeter, $P$ is
  \begin{equation}
    P=2l +2w\label{eq:SquareConst}
  \end{equation}
  and is constrained to be a constant (fixed) value so
  equation~\ref{eq:SquareConst} is the constraint. 

  At this point we can follow either of  two equivalent procedures.
  \begin{description}
  \item[Leibniz's Approach:] 
    \begin{enumerate}
    \item Compute $\dx{A} = l\dx{w}+w\dx{l}$, from the objective
      function.
    \item Differentiate the constraint, giving the differential form of
      the constraint: $2\dx{l}+2\dx{w}=0$ 
    \item Set $\dx{A}=0$ and solve using the constraint itself and its
      differential form as needed.
    \end{enumerate}
    In our opinion this is how Leibniz intended for his Calculus to be
    used so we tend to favor it.
  \item[The Modern Approach:] 
    \begin{enumerate}
    \item Solve the constraint for either $w$ or $l$ in terms of the
      other. For example, $l=\frac12(P-2w).$
    \item Substitute this value into the objective function:
      $A=\frac12(P-2w)\cdot w$. This gives the area as a function of
      the width ($w$) alone, so we can now think of $A(w)$ as an
      abstract function and find its maximum by setting
      $\dfdx{A}{w}=0$ and solving.
    \end{enumerate}
    The modern approach emphasizes the use of functions and
    derivatives rather than curves and differentials.
  \end{description}
  These two approaches are completely equivalent so you can use
  whichever feels most comfortable to you. We tend to
  favor the first approach.
  
  \begin{embeddedproblem}
    Solve this problem using both methods to see that they are
    equivalent.
  \end{embeddedproblem}
\end{myexample}

\begin{ProficiencyDrill-1line}
  \begin{enumerate}[label={\bf{}(\alph*)}]
  \item Find the approximate coordinates of the point on the curve
    $y=e^x$ closest to the origin.
  \item Verify that this is also where the line connecting this point
    to the origin is perpendicular to the curve.
  \end{enumerate}
\end{ProficiencyDrill-1line}

\begin{ProficiencyDrill}
  Suppose a point is moving so that its $x-y$ coordinates are given
  by: $
  \begin{pmatrix}
      x(t)=t-1\\
      y(t)=e^t
    \end{pmatrix}$.
    \begin{enumerate}[label={\bf{}(\alph*)}]
      \item Find the value of $t$ where the point is closest to the
        origin.
      \item Find the $x-y$ coordinate so of the point's location when
        it is closest to the origin.
    \end{enumerate}
  \end{ProficiencyDrill}
  

\begin{ProficiencyDrill-1line}
  \label{problem:RectInscEll}
    \begin{enumerate}[label={\bf{}(\alph*)}]
    \item     Of all the rectangles that can be inscribed in the unit circle:\\
      \centerline{\includegraphics*[height=1in,width=1in]{../Figures/InscribedSquare1}}
      show that the square has the largest possible area. (Warning: Do
      not assume that we have drawn the best diagram possible for
      you.)
    \item     Of all the rectangles that can be inscribed in the
      ellipse $x^2+\frac{y^2}{9}=1$ show that the one with the maximum
      area has length three times longer than its width.
    \item     Of all the rectangles that can be inscribed in the
      ellipse $\frac{x^2}{16}+y^2=1$ show that the one with the maximum
      area has length four times longer than its width.
    \item Of all the rectangles that can be inscribed in the ellipse
      $\frac{x^2}{a}+\frac{y^2}{b}=1$ show that for the one with the
      maximum area the ratio of length to width is
      $\sqrt{\frac{b}{a}}$.
    \end{enumerate}
\end{ProficiencyDrill-1line}


We have seen (in several different ways) that out of all rectangles
with fixed perimeter, the square (an equilateral rectangle) has the
largest area. Is this also true of triangles? More precisely, does an
equilateral triangle have the largest area out of all triangles with a
fixed perimeter?  This may sound like a simpler problem but it is
actually quite a bit harder.

  With a rectangle the parallel sides varied together so we only had
  two parameters to deal with. For the triangle all three sides vary
  independently which gives us three independent parameters.

  Rather than addressing the equilateral triangle problem directly
let's try a related, but simpler, problem. This is a standard
problem-solving ploy. We replace a hard problem with an ``adjacent''
problem that is (we hope) more tractable. The hope is that the simpler
problem will give us insight into the harder problem.

For example, suppose we hold one of the sides of our triangle
fixed. Now we only have two variables again.

\begin{wrapfigure}[]{r}{2in}
  \captionsetup{labelformat=empty}
  \centerline{\includegraphics*[height=1in,width=2in]{../Figures/Optimization3}}
\label{fig:}
\end{wrapfigure}
So the problem reduces to this one: given all triangles with a fixed
base and a fixed perimeter, does an isosceles triangle have the
largest area?  The sketch at the right has all of  the pertinent
information labeled.

Since the area $A=\frac12bh$ and $b$ is being held fixed, then
maximizing $A$ amounts to maximizing $h$, so $h$ is the natural
objective function for this problem.  For reasons that should be clear
we will replace our objective function $h$ by $H=h^2$ and work on
maximizing this.  The twist in this problem is that we actually have
two different formulas for $H$.
\begin{embeddedproblem-enumerate}
  \begin{enumerate}[label={\bf{}(\alph*)}]
    \item Express $H$ as a function of $x$ and
      $z$.  Differentiate this and set $\dx{H}=0$ to obtain an
      equation relating $\dx{y}$ and $\dx{z}$. 
    \item Express $H$ as a function of $y$ and
      $w$. Differentiate this and set $\dx{H}=0$ to obtain an equation
      relating $\dx{y}$ and $\dx{w}$.
    \item Use the fact that $b$ is constant to find an equation
      relating $\dx{z}$ and $\dx{w}$.   Use the fact that the
      perimeter $P$ is constant to find an equation relating $\dx{x}$
      and $\dx{y}$.
    \item Use the results from parts (a), (b), (c) to show that $H$ is
      maximized when
      $$
      \frac{-x}{z}=\dfdx{z}{x}=\frac{-y}{w}
      $$
      and use this to conclude that the base angles of the triangle
      must be congruent and the triangle with the largest area must be
      isosceles.
  \end{enumerate}
\end{embeddedproblem-enumerate}

That was the ``adjacent'' problem. How can we use this to show that
the maximal area of a triangle with a fixed perimeter must be
equilaterl.  We now have all of the pieces of an argument. But it
probably isn't yet clear how to put them together.

We will require some indirect reasoning as described by Sherlock
Holmes in the novel \pubtitle{The Sign of the Four}: ``How often have
I said to you that when we have eliminated the impossible, whatever
remains, \emph{however improbable}, must be the truth?''

For our problem we will determine that if a triangle is not
equilateral, then a triangle with the same perimeter can be made with
a larger area.  Assuming there is a largest triangle, then the only
possibility left is that it must be equilateral.

\begin{embeddedproblem}
  \label{problem:MaximalTriangle}
  Suppose we have a triangle with sides $a$, $b$, and $c$ where at
  least two sides, say $a$ and $c$, are not equal. We want to show
  \foreign{\'{a} la} Holmes that it is impossible for such a triangle
  to have maximal area.
  \begin{enumerate}[label={\bf{}(\alph*)}]
  \item If we hold $b$ fixed explain how the previous problem implies
    that there must be another triangle with the same perimeter and a
    larger area. 
  \item Explain how this implies that the maximal triangle must be
    equilateral.
  \end{enumerate}
\end{embeddedproblem}


% \begin{enumerate}[label={\bf{}(\alph*)}]
%     \item Try attacking the largest triangle problem directly, in the
%       same way we attacked the largest rectangle problem. The
%       following diagram might help:\\
%       \centerline{\includegraphics*[height=2in,width=4in]{../Figures/Optimization3}}
%       From the diagram we see that the area of the triangle is
%       $$
%       A=\frac12ca\sin(\alpha)=\frac12cb\sin(\beta).
%       $$
%       So that is the objective function. What constraints do we have?
%       Can you solve this problem?
%       If you can't, don't worry about it. This was hard. There are
%       five diffent variables involved ($a, b, c, \alpha,$ and $\beta$)
%       in play here. Clearly it will help to reduce the number of
%       variables.
%     \item Try this instead. An ancient Greek mathematician named Heron
%       (or Hero) of Alexandria ($10-75$ AD) showed that the area of a triangle can
%       be determined from the lengths of the sides of the triangle via
%       the formula:
%       $$
%       A=\sqrt{\frac{P}{2}\left(\frac{P}{2}-a\right)\left(\frac{P}{2}-b\right)\left(\frac{P}{2}-c\right)},$$
%       where $P=a+b+c.$

%       If you use Heron's Formula then we only have three variables in
%       play but now there is the square root involved. Does that
%       complicate things?

%       No, not really. We say in examples~\ref{example:LineDist}
%       and~\ref{example:CircleDist} that optimizing a quantity, and
%       optimizing its square are essentially the same problem so we
%       could simply let
%       $$
%       O(x) = A^2 =
%       \frac{P}{2}\left(\frac{P}{2}-a\right)\left(\frac{P}{2}-b\right)\left(\frac{P}{2}-c\right)
%       $$
%       be the objective function.

%       What are the constraints in this case? Can you solve the problem
%       now?
%     \item Even with Heron's Formula this is still a formidable
%       problem. Is there something we can do to simplify things
%       further? We've had some success with problems with two
%       variables. Is there some way to reduce the number of variables
%       by one more.
%     \item When faced with a hard problem (which this seems to be) a
%       good strategy is to look for a similar, related problem to solve
%       to see if that sheds any light. In practice finding a good
%       ``adjacent'' problem can be difficult and time-consuming all by
%       itself.

%       Since we can't spend time chasing down bad ideas we'll simply
%       tell you that the adjacent problem you want to look at is the
%       one where only two sides are free to change. Think about
%       it. We've had success with two variable problems and we've
%       managed to whittle this one down to three variables when we use
%       Heron's Formula. We can reduce it to two variables by simply
%       holding one side of the triangle fixed.

%       Of course this isn't the same problem but maybe, if we can solve
%       it, we'll see how to extend it to our problem.

%       So the problem is: Of all triangles with sides $a$, $b$, and
%       $c$, where $a$ is constant, and the perimeter $P=a+b+c$ is also
%       constant, show that the one with the largest area is
%       isoceles. (That is, $b=c$.) \hint{Use the square of Heron's
%         Formula for the objective function.}
%     \item Ok. Can you see how to address our problem now?

%       No?

%       Let's try this: Suppose the triangle with maximum area is
%       \alert{not} equilateral. That means the two of the sides, say
%       $a$ and $b$ are unequal in length. Whatever the length of $c$
%       is, consider it fixed. Now we have a maximum area triangle with a fixed base
%       ($c$) but it is not isoceles because $a$ and $b$ are unequal.

%       What do you conclude from this?
%     \end{enumerate}
% \end{embeddedproblem}

\section{Reflections, Refraction, and Rainbows}
\label{sec:refl-refr-rainb}
Because they are called ``Optimization Problems'' it is natural to
assume that such problems are always about finding the maximum
or minimum value of something. To be sure, many are. But not all.

For instance, consider Problem~\ref{problem:MaximalTriangle} above. It
did not ask you to find the the maximum possible area of the triangle
(although that would be easy to compute once you've solve the
problem). Instead it asked you to describe a property of the the
maximal triangle.  This illustrates that optimization problems come in
two flavors:
\begin{enumerate}
\item Optimize some objective function, or
\item Show that some property holds when some objective function is
  optimized.
\end{enumerate}

Another example of the second type is the problem of reflecting light
we looked at in chapter~\ref{chapter:ScienceBeforeCalc}. Recall that we
showed that light bounces off a mirror so that the angle of incidence
equals the angle of reflection.  We did this by reasoning  geometrically
from the figure below:
  \centerline{\includegraphics*[height=1.4in,width=3.5in]{../Figures/ReflectingMirror4}}
  to show that the minimal path must satisfy the condition $\angle\,
  1=\angle\, 2$.
  
  The geometric argument was actually pretty easy to understand once
  we found the right trick. But that's the difficulty. The geometric
  argument required the very counterintuitive ``trick'' off pretending
  that the light passed through the mirror instead of reflecting off
  of it!

Calculus provides a more systematic approach, one that does not
require a ``stroke of genius'' level insight to solve each new 
problem. We'll redo this problem, this time using the Calculus tools
we've learned.

% Calculus gives us our systematic way to look at this.  Let's solve
% this problem again using the techniques from the last section.
\subsection{Reflection}\label{subsec:reflection}

First, we need an objective function.  Drawing a labeled picture is a
good start, so we modify the sketch above to suit the needs of
Calculus. In particular, the relevant quantities need names.
  \begin{figure}[h]
    \centering
    \includegraphics*[height=2in,width=5in]{../Figures/ReflectingMirror5}  
    \caption{\textcolor{blue}{Reflecting Light}}
    \label{fig:RefMir}
  \end{figure}

    We assume that $a, b$, and $c$ are constants while $x, y, z$, and
    $w$ are variables.

    We need to argue that when the path of the light from $A$ to $R$
    to $B$ is shortest, then $\alpha=\beta$. Since the length of that
    path is (from our sketch) $z+w$ the most natural objective
    function is: 

% Our argument consisted essentially of assuming that a ray of light
% will take the shortest possible path and using this to show that on
% this path the angles are equal. This is an optimization problem
% but notice that in this problem we are not actually interested in the
% minimum length traveled by the lightbeam. (This would be the sum $z+w$
% in our diagram.)

% of a different kind. The quantity being optimized
% here is clearly the length of the path followed by the light beam. But
% we weren't actually interested in the value of that length. We wanted
% to show that \emph{if} the light followed the shortest possible path
% then the angles of incidence and reflection had to be the same.

% The geometric reasoning we used is simple enough to understand when it is explained but a good deal
% of ingenuity would have been required to think of it in the first place. This is the sort of thing
% that prompted Leibniz to remark, ``.\ .\ . very learned [people] have sought in many devious
% ways what someone versed in this calculus can accomplish in these
% lines as by magic.''



% Thus the obvious place to look for extrema are those points where
% $\eval{\dfdx{y}{x}}{x}{a}=0$.  We'll call such a value a
% \term{possible transitional value} for reasons that will become apparent later.

% But Fermat's Theorem does \emph{not} say that if
% $\eval{\dfdx{y}{x}}{x}{a}=0$ then the point $(a,b)$ must be an
% extremum. In fact, this is not true. We will have more to say
% about this later. For now we know what the first step in the
% search for extrema must be: Find the derivative and set it equal
% to zero.  Let's apply this approach to the reflection problem.

% Strictly speaking, Fermat's Theorem also imposes the condition
% that $f$ be differentiable at $x=a$. This is a real condition
% and we will have more to say about it later. For now we simply
% observe, via this problem, that extrema can also occur where the
% function cannot be differentiated.

  \begin{equation}
    \label{eq:RefMir1} L=z+w
  \end{equation}
  and we know from Fermat's Theorem that $L$ will be minimum when
  $\dx{L}=0$.

   We have the following  constraints
   \begin{align}
    \label{eq:RefMir2}  a^2+x^2&=z^2, \text{ and}\\
    \label{eq:RefMir3} b^2+y^2&=w^2.\\
    \label{eq:RefMir4} x+y&=c,\\
  \end{align}


  Differentiating equation~\ref{eq:RefMir1} gives $\dx{L} = \dx{z}+\dx{w}$.

  Differentiating equation~\ref{eq:RefMir2} gives
  $\dx{z}=\frac{x}{z}\dx{x}$, so
  $$
  \dx{L} = \frac{x}{z}\dx{x}+\dx{w}.
  $$

  Differentiating equation~\ref{eq:RefMir3}
  gives$\dx{w}=\frac{y}{w}\dx{y}$, so
  $$
  \dx{L} = \frac{x}{z}\dx{x}+\frac{y}{w}\dx{y}.
  $$

  And finally differentiating equation~\ref{eq:RefMir4} gives $\dx{y} =
  -\dx{x}$, so that
  $$
  \dx{L} = \frac{x}{z}\dx{x}-\frac{y}{w}\dx{x},
  $$
  or
  $$
  \dfdx{L}{x} = \frac{x}{z}-\frac{y}{w}.
  $$


  Now suppose that the path from $A$ to $B$ shown in the sketch is the
  shortest possible path, i.e., that $L$ is minimal. Then by Fermat's
  Theorem it must be that $\dfdx{L}{x}=0$, and therefore that
  $$
  \frac{x}{z}=\frac{y}{w}.
  $$

  \begin{ProficiencyDrill}
    Use the fact that the shortest path occurs when
    $\frac{x}{z}=\frac{y}{w}$ to conclude that $\alpha=\beta$.
  \end{ProficiencyDrill}


  We will readily agree that this Calculus based argument is not as
  nice as the geometric argument we used in
  chapter~\ref{chapter:ScienceBeforeCalc}. In particular the geometric
  argument required a new sketch at every step so we could easily
  visualize what was happening.

  Our Calculus based argument wasn't at all visual and it involved a
  lot of computation. But, believe it or not, this is a strength not a
  weakness.  It is a strength because none of the computations were
  particularly hard. They were just tedious and detailed. But -- and
  this is the point -- it also gave us a clear map of what
  computations needed to be done.  We needed to
  \begin{enumerate}
  \item Compute  $\dfdx{L}{x}$, and
  \item  Set $\dfdx{L}{x}$ equal to zero.
  \end{enumerate}
  The rest was
  algebra which is straightforward (though not necessarily easy). This
  is why you need to be skilled at differentiation. If any of the
  differentials we obtained above had been wrong it wouldn't matter
  how good our algebra skills are.

  In this problem the ultimate goal, showing that $\alpha=\beta$,
  emerged almost immediately so we stopped there. This will not always
  be the case.  Suppose for a moment that we hadn't noticed that
  $\frac{x}{z}=\cos(\alpha)$ and $\frac{y}{w}=\cos(\beta)$. We could
  still solve this problem by finding the lengths $x$ and $y=c-x$ and
  from there the equivalence of $\alpha$ and $\beta$ follows.

\begin{embeddedproblem}
  In figure~\ref{fig:RefMir} suppose that $a=5$, $b=3$, and $c=10$.
  \begin{enumerate}[label={\bf{}(\alph*)}]
  \item Show that
    $\dfdx{L}{x}=\frac{x}{\sqrt{25+x^2}}-\frac{10-x}{\sqrt{9-(10-x)^2}}$.
  \item Use Newton's Method to show that
    $L\left(\frac{25}{4}\right)\approx0$.
  \item Use the result of part (b) to show that $\alpha=\beta$.
  \end{enumerate}
  This solution is not pretty or elegant, but it is
  effective. Moreover no cleverness is needed, just computational
  competence. 
\end{embeddedproblem}


\subsection{Refraction: Snell's Law}\label{subsec:refraction}
  We first addressed Snell's Law in Section~\ref{sec:snells-law-refr}
  of Chapter~\ref{chapter:ScienceBeforeCalc}. Recall that in
  PIC~\ref{PIC:FermatSnell} we had the following sketch \\
  \centerline{\includegraphics*[height=2.2in,width=3.7in]{../Figures/Refraction4}}
  which represents the passage of light from $A$ to $C$ to $B$. We
  assumed that velocity of light from $A$ to $C$ is $v_1$ and from $C$
  to $B$ is $v_2$. In that case Snell's Law says that:
  $$
  \frac{\sin(\alpha)}{v_1}=  \frac{\sin(\beta)}{v_2}.
  $$

  We did not derive Snell's Law geometrically in
  Chapter~\ref{chapter:ScienceBeforeCalc} because, like reflection in
  the previous section, as small stroke of genius is needed to set up
  the argument.

  Willebrord Snell  first  demonstrated Snell's Law (that's why
  it's named after him) in $1621$ well before the invention of Calculus but his
  demonstration wasn't  published until $1703$ by  Leibniz's
  teacher, Christiaan Huygens. We did not show you a pre-Calculus 
  proof of Snell's Law in
  section~\ref{sec:snells-law-refr} for the simple reason that it is \emph{hard}.
  Instead we asked you in problem~\ref{PIC:FermatSnell}
  to get started on the problem by finding an expression for the time,
  $T$ needed for light to travel along the path from $A$ to $C$ to $B$
  in terms of the constants $a, b, c, v_1, v_2$ and the variable $x$.

  \begin{ProficiencyDrill}
    \label{PD:SnellMinFunc}
    Show that the solution of  problem~\ref{PIC:FermatSnell} is:
    $$
    T(x)=\frac{\sqrt{a^2+x^2}}{v_1}+\frac{\sqrt{b^2+(c-x)^2}}{v_2}.
    $$
  \end{ProficiencyDrill}

Having obtained $T(x)$ we would seem to be all set to derive Snell's
Law by solving $\dfdx{T}{x}=0$ for $x$.

 % and it is not clear how you get
  % $\frac{\sin(\alpha)}{v_1}=\frac{\sin(\beta)}{v_2}$ out of that.
  % Instead, let's look at our objective function and constraints and go
  % from there

  \begin{embeddedproblem}
   Try it. That is, compute the derivative of $T$ from Proficiency
   Drill\ref{PD:SnellMinFunc}. 

   The differentiations necessary to compute $\dfdx{T}{x}$ are
   somewhat challenging but if you proceed carefully they are not
   insurmountable. The real issue is solving $\dfdx{T}{x}=0$ for
   $x$. Show that when you set $\dfdx{T}{x}=0$ you get the following
   equation to solve:
   $$
   (v_2^2-v_1^2)x^4-2(v_2^2-v_1^2c)x^3+\left(v_2^2(b^2+c^2)+v_1^2(a^2+c^2)\right)x^2+v_1^2a^2c(2-c)
   =0.
   $$

   The next step would be to solve this equation for $x$ (it's as bad as
   it looks) and then show that at this optimal value for $x$ Snell's
   Law holds.

   We won't ask you do do that.  
 \end{embeddedproblem}
 
    However, recall that in this section and in this problem, we
    weren't trying to actually find a minimum value of $T(x)$. That
    was just a step along the way to our real goal: Finding a
    condition 
    that holds when we have a optimum.   The
    condition of course, is Snell's Law.

    With this in mind, let's go back to
    our diagram and not worry about trying to get everything in terms
    of one variable.  Instead, we will use as many variables as we
    need and write down constraints on them.

    \begin{embeddedproblem-enumerate}
    \begin{enumerate}
    \item We already know that the total time traveled from point $A$  to $C$ to $B$ is given by
      $$
      T(x)=\frac{z}{v_1}+\frac{w}{v_2}.
      $$
      Write down all of the constraints involving  the variables $x$,
      $y$, $z$, and $w$, and the constants $a$, $b$, and $c$.  (As before, we'll ignore $\alpha$ and $\beta$ for now.)
    \item Show that $T$ is minimized when
      $\frac{1}{v_1}\frac{x}{z}=\frac{1}{v_2}\frac{y}{w}$ and use this
      to derive Snell's Law of refraction:
      $$
      \frac{\sin(\alpha)}{v_1}=\frac{\sin(\beta)}{v_2}.
      $$
 (This should look a lot like what we
      did in example~\ref{example:ReflectingMirror}.)
    \end{enumerate}
    \end{embeddedproblem-enumerate}

  % \subsection{Rainbows}\label{subsec:rainbows}

Of course light reflects off of everything, not just flat mirrors.  A
general discussion of that phenomenon is more properly done in a
course on classical physics.  We will limit ourselves to the geometry
involved when a light beam strikes a spherical mirror, like the
surface of a raindrop. Since a raindrop  is not a perfect mirror some
of the beam passes through the mirror and refracts, while some of it
reflects back off of the surface.

\begin{wrapfigure}[]{r}{2in}
  \captionsetup{labelformat=empty}
  \centerline{\includegraphics*[height=1in,width=2in]{../Figures/Refraction10}}
  \label{fig:Refraction10}
\end{wrapfigure}
Happily nothing changes much. We need only observe that light only
interacts with the mirror locally, where it actually strikes  the
mirror. And since we know that at very small scales any curve is
indistinguishable from it's tangent line it should be clear that when
a light beam strikes a curved mirror it reflects or refracts exactly
as if it were striking a flat mirror tangent to the curve at the point
of interaction as seen at the right. This is again the Principle of Local
Linearity  

\begin{wrapfigure}[]{l}{2.1in}
  \captionsetup{labelformat=empty}
  \centerline{\includegraphics*[height=.75in,width=2.1in]{../Figures/DoubleRainbow6}}
  \label{fig:DoubleRainbow6}
\end{wrapfigure}
The ordering of the colors of a rainbow, the shape of a
rainbow, its position relative to the viewer, and the simple fact that
rainbows exist to be seen are a result of the reflective and
refractive properties of light. Even better for our purposes
optimization is needed to answer them.

Newton determined experimentally that white light was a mixture of
other colors. It was well known that the time that a beam of light
passing through a prism would emerge with different colors bent
(refracted) at different angles as seen above on the left. A prevailing
theory of the time was that the prism had altered the light in some
way that had added the colors. Thus it was believed that the color
bands that emerged from the prism were a property of the prism, not of
the light.

\begin{wrapfigure}[]{r}{4in}
  \captionsetup{labelformat=empty}
  \label{fig:DoubleRainbow9}
  \centerline{\includegraphics*[height=1.9in,width=4in]{../Figures/DoubleRainbow9}}
  \caption{The double rainbow on the left was captured by Elizabeth
    Eschel (a friend of one of the authors) of Tsfat/Safed, Israel in
    2018 while on her way to
    work.}\end{wrapfigure}
Newton debunked that theory by passing \emph{only} the red band
through a second prism. He reasoned that if the first prism had
somehow altered the light then the second would alter it again. But
when only red light emerged from the second prism, refracted but
otherwise unchanged, he concluded that the color was an inherent
property of the light, not the prism.

% wrote a seminal work on optics and .  This can be seen as light passes
% through a prism or in a rainbow.  Surely you heard the mnemonic
% \alert{ROY G BIV} for the colors of a rainbow.  Which color is on top,
% red or violet?  Before you go further, take a guess.  % \\
Look closely at the three rainbows in the images at the right above. They
share certain properties that we would like to understand. These are:
\begin{enumerate}
\item They are bright. At least they are brighter than the surrounding
  viewscape. Indeed it is only the brightness of the light coming from
  the apparent location of the rainbow that makes it visible since a
  rainbow is not a physical object.
\item They have a particular shape, but what is it? Are they
  circular? Parabolic? Elliptical? Some other, possibly unnamed
  shape?
\item Two of the rainbows shown are double rainbows. Notice that the
  secondary rainbow is dimmer\footnote{That's actually what makes it
    the secondary.}. Also the colors are reversed from the primary to
  the secondary. We'd like to understand that.
\end{enumerate}
Each of the properties listed above can be couched as an optimization
problem. We will address them one at a time.

To keep things as simple as possible we will assume that water
droplets suspended in the atmosphere assume a perfectly spherical
shape. This is close enough to reality that the errors in our
computations will not be significant.


Also, we will analyze this as if everything were two-dimensional. That
is we will draw the light beam crossing into a circular raindrop
rather than a spherical raindrop. Again this is just to keep things
simple. The analysis of the three dimensional problem  yields
similar results but is much harder algebraically.

\begin{wrapfigure}[]{r}{1.7in}
  \captionsetup{labelformat=empty}
  \centerline{\includegraphics*[height=1.7in,width=1.7in]{../Figures/Refraction6}}
  \label{fig:Refraction6}
\end{wrapfigure}
As we've seen when light refracts across the plane dividing air and
water it bends toward the normal (perpendicular). The physical reason
for this is that the velocity of light in air is faster than its
velocity in water. If it transitions from a slower medium (like water)
to a faster medium (like air) then it bends away from the normal. Thus
in the diagram at the right we can deduce that the velocity is lower
in the region below the boundary (in black) whichever direction the
light ray (in red) is moving.



% Of course the sketch above doesn't capture the curvature of the rain
% droplet.  Our sketch should look more like the figure at the left,
% where the light beam is crossing the black, circular boundary of the
% droplet. Does Snell's Law still apply when the boundary is curved?

% Of course it does. The light beam doesn't care what shape the boundary
% has away from the point of entry. It only interacts with the boundary
% locally, and locally any curve is indistiquishable from it's tangent
% line so the light beam will refract as it crosses into the droplet
% just as if it were crossing the line tangent to the circle at the
% point of intersection. As we saw in\notetoself{We need to add this
% problem somewhere before this, if we don't already have it: Show that the line tangent to a
% circle is orthogonal to the radius of the circle.  -- Sep 3, 2020}
% Problem~\#\ref{} the line tangent to a circle at a point is orthogonal
% to the radius that passes through the same point. The radius and
% tangent are the blue, dotted lines in the diagram at the left.
But the physics involved is not our concern. All we need is  the
knowledge that Snell's Law is in play in the diagram at the left.
$\frac{\sin(\alpha)}{v_1}=\frac{\sin(\beta)}{v_2}.$
The green line in the diagram is the path the light beam would have
followed had it \emph{not} refracted at the boundary.

\begin{ProficiencyDrill}
  For later reference notice that the angle between the red and green
  lines is $\alpha-\beta$ as indicated. \hint{This is an elementary
    result from geometry. Don't make it hard.}
\end{ProficiencyDrill}
We now have everything we need to trace the path of a beam of light
into and back out of a water droplet.
In the sketch below the red arrows indicate  the path the light
beam actually follows, $O$ is
the center of the droplet, and the line segments $OA,
OB,$ and $OC$ are all radii of the droplet.

% \begin{figure}[h]
  
  \centerline{\includegraphics*[height=2.5in,width=4in]{../Figures/Refraction7}}
  
%   \caption{The path of a typical light beam as it moves through a
%     spherical droplet.}
%   \label{fig:DoublePath}
% \end{figure}
 
In order to address the question of the rainbow's brightness -- why it
is even visible in the first place -- we want to find the total
clockwise rotation of the direction of the beam from the time it
enters the raindrop to the time it exits. 

\begin{wrapfigure}[]{r}{1.5in}
  \captionsetup{labelformat=empty}
  \centerline{\includegraphics*[height=.6in,width=1.5in]{../Figures/Refraction8}}
  \label{fig:Refraction8}
\end{wrapfigure}
That is, if we place the
starting and ending beams head-to-tail as indicated in the sketch at the right we
want to find the angle $D(\alpha)$. This is called the
\term{deviation angle} and for a given entry angle $\alpha$ it tells
us what the exit angle will be. We are interested in finding the value
of $\alpha$ that \emph{minimizes} $D(\alpha)$.

Here's why. Suppose $D(\alpha)$ is minimum for some angle, say
$\alpha=\phi$, and that $D(\phi)=\psi$ is the minimum. Then all of the
light beams entering the droplet at an angle \emph{near} $\phi$ will
exit \emph{near} $\psi$. It is generally true that things which start
out close together usually end up close together. Functions with this
property are called \term{continuous functions}.

\begin{wrapfigure}[]{l}{3.5in}
  \captionsetup{labelformat=empty}
  \centerline{\includegraphics*[height=1.5in,width=3.5in]{../Figures/DoubleRainbow3}}
  \label{fig:DoubleRainbow3}
\end{wrapfigure}
But the nature of the function
$D(\alpha)$ guarantees that there is a large interval of entry
angles which will all have approximately the same exit angle.  The
easiest way to see this is to look at the the shape of the graph of
$D(\alpha)$ below. Notice that the
black box we drew around then minimum point (denoted\aside{For now you
  will have to trust that we have represented $D(\alpha)$
  accurately. We will shortly be deriving a formula for $D(\alpha)$
  and you can verify the accuracy of this graph yourself.}
$(\gamma,\psi)$ in the sketch.). is much wider than it is
tall.
% This says that all of the the light beams entering at an angle
% between about $0.8$ radians and $1.2$ radians will exit at about $2.4$
% radians. That is the light will be concentrated (brighter) at a
% deviation angle of about $2.4$ radians (about $138^\circ$. 

% We will soon see that the deviation angle
% depends \emph{only} on the angle $\alpha$ as we've
% indicated. 
From the diagram it is clear that for entry angles between a little
more than $0.8$ radians ($\approx45.8^\circ$) and a little more than
$1.2$ radians ($\approx68.8^\circ$) the exit angles will all be very
close to $2.4\approx138^\circ$ radians. Outside of the interval
$(.8, 1.2)$ the exit angles grow very quickly.

This says that the light beams coming out of a droplet at an angle of
$\psi\approx2.4$ radians are a concentration of the beams that entered
between about $.8$ and $1.2$ radians. So light
exiting at about $2.4$ radians will be brighter than the surrounding
light. This is what makes a rainbow visible in the first place.

Now of course, we need to do the calculations to confirm all of
this. To find a formula for the deviation angle, $D(\alpha)$ we
follow the path of a typical light beam as it refracts and reflects
its way into and then back out of a typical droplet.

In Figure~\ref{fig:DoublePath} we see that the beam (in red) makes a
turn of $\alpha-\beta$ at the point $A$ and another turn of the
same magnitude at point $C$ so
$$
D(\alpha) =
2(\alpha-\beta) + \{\text{the magnitude of the rotation at point
  $B$}\}.
$$

\begin{wrapfigure}[]{r}{2in}
  \captionsetup{labelformat=empty}
  \centerline{\includegraphics*[height=1.8in,width=2in]{../Figures/Refraction9}}
  \label{fig:Refraction9}
\end{wrapfigure}
To find the magnitude of the rotation at point $B$ consider the sketch at
the right. Clearly at point $B$ the beam turns by the angle labeled
$\theta$.

\begin{ProficiencyDrill}
  \begin{enumerate}[label={\bf{}(\alph*)}]
  \item   Use Geometry to show that $\alpha=\angle AOB$.
  \item Show that $\alpha=\pi-2\beta$. \hint{The interior angles
      of triangle $OAB$ must sum to $\pi$. What is $\beta+\theta$
      equal to?}
  \end{enumerate}
\end{ProficiencyDrill}
Therefore
\begin{align*}
  D(\alpha) &= 2(\alpha-\beta) + \{\text{the magnitude of the
              rotation at point  $B$}\}\\
            &= 2(\alpha-\beta) + \alpha\\
            &= 2(\alpha-\beta) + \pi-2\beta, \text{ or}\\
            &= \pi+2\alpha-4\beta.
\end{align*}

Notice that $\alpha$ and $\beta$ are not independent, they are governed by Snell's
Law which we will rewrite as
\begin{equation}
  \sin(\alpha)=\frac{v_{\alpha}}{v_{\beta}}\sin(\beta)=k\sin(\beta).\label{eq:IndRef}
\end{equation}
The constant $k=\frac{v_{\alpha}}{v_{\beta}}$ is called  the ``index of refraction'' (for
water).

\begin{ProficiencyDrill}
  Show that $\beta=\inverse\ \lt{}sin\left(\frac{\sin(\alpha)}{k}\right)$.
\end{ProficiencyDrill}
Knowing $\beta$ we can now express the deviation angle,
$D(\alpha)$ as a function of $\alpha$ alone:
\begin{align*}
  D(\alpha) &= \pi+2\alpha-4\beta \\
  D(\alpha) &= \pi+2\alpha-4\inverse\sin\left(\frac{\sin(\alpha)}{k}\right).
\end{align*}

\begin{embeddedproblem}
  For visible light  the index of refraction
  is $k\approx\frac43$. Mathematically, the problem becomes minimizing
  $D(\alpha)=\pi+2\alpha-4\beta$ subject to the constraint
  $\sin(\alpha)=k\sin(\beta)$, $ 0\le\alpha\lt{}\frac{\pi}{2} $.

  \begin{enumerate}[label={\bf{}(\alph*)}]
  \item Graph $D(\alpha)$ with $k=\frac43$ and confirm that the graph
    we showed you earlier is correct.
  \item Use Newton's Method to show that the coordinates,
    $(\gamma,\psi),$
    of the minimum point on the graph are $\gamma\approx1.04\approx59.59^\circ$ and
    $\psi\approx2.41\approx138.1^\circ$ respectively.
  \end{enumerate}
\end{embeddedproblem}
% ****************** Pick up here *************************************

% Graphing $D(\alpha)$ for this value of  $k$ and
% $0\le{}\alpha{}\lt{}\frac{\pi{}}{2}$, we have.
% Notice from the graph that the minimum of this curve appears to be
% around the point $(1.03, 2.41)$.  This says that the minimum value of
% $D(\alpha)$  is about $2.41\text{ radians }\approx 138^\circ$  and it
% occurs when $\alpha\approx1.03\text{ radians }\approx59.4^\circ$.
% As we can see from our sketch the  minimum 
% value of $\alpha$ near  $1.03$  the amount of deviation is roughly
% the same.

% \vskip.5in{}
% ***************************************************************\\
% As you can see in the diagram below, the ray of light entering the
% raindrop at an angle of radian measure $\alpha$ is refracted to an
% angle of radian measure $\beta$.  This is governed by Snell's Law.\\
% % \begin{wrapfigure}[]{l}{3in}
% %   \captionsetup{labelformat=empty}
% %   \centerline{\includegraphics*[height=1.75in,width=3.25in]{../Figures/DoubleRainbow2}}
% %   \label{fig:DoubleRainbow2}
% % \end{wrapfigure}
% $\frac{\sin(\alpha)}{v_\alpha}=\frac{\sin(\beta)}{v_\beta}$, where
% $v_\alpha$ and $v_\beta$ are the velocities of light in air and water,
% respectively.  Actually, only a portion is refracted as the other
% portion is reflected off the exterior of the raindrop.  Once inside
% the raindrop, it reflects off of the back of the raindrop (again only
% a portion) so that the angle of incidence equals the angle of
% reflection.  A portion of the light then exits the raindrop governed
% by Snell's Law again.  The total amount of (clockwise) rotation in
% this process is called the deviation angle, $D(\alpha)$.%

% \begin{embeddedproblem}
%   Show that $D(\alpha)=\pi-4\beta+2\alpha.$
% \end{embeddedproblem}


\begin{wrapfigure}[]{r}{2.5in}
  \captionsetup{labelformat=empty}
  \includegraphics*[height=.8in,width=2.5in]{../Figures/DoubleRainbow5}
  \label{fig:DoubleRainbow5}
\end{wrapfigure}
This says that light rays entering the raindrop at approximately
$59^\circ$ have the highest concentration when exiting the raindrop.
At other values of $\alpha$ the dispersion of light is greater.  This
provides the angle that the rainbow makes with an observer.  This
supplementary angle to $D(\alpha)$ is called the rainbow angle,
$$
R(\alpha) =\pi-D(\alpha)=\pi-(\pi+2\alpha-4\beta)=4\beta-2\alpha\text{
  radians.}
$$
For white light, we have a rainbow angle of
$$
R(1.03)\approx\pi-2.41\approx0.734\text{ radians}\approx42^\circ.
$$

To obtain the spread of colors, note that the index of refraction
changes with the color of light.  We said earlier that the index of
refraction for visible light is approximately $\frac43$. This is true,
but it actually varies with the frequency (color) of the light. The
red light refracts at one angle and the orange at a slightly different
angle so they appear in the rainbow adjacent to each other. This
variation is continuous throughout the spectrum which is what gives us
the rainbow.


\begin{embeddedproblem}
  We could use Newton's Method like we did earlier to find the
  minimum of $D(\alpha)$ for various values of $k$ but it is a
  little easier to do it this way:
     \begin{enumerate}[label={\bf{}(\alph*)}]
  \item Show that $D(\alpha)$ is minimized when % $\cos(\alpha)=\frac{k}{2}\cos(\beta)$.
  \item Use the equation obtained in part (a) along with the
    constraint $\sin(\alpha)=k\sin(\beta)$ to show that $D(\alpha)$
    is minimized when $\sin(\beta)=\sqrt{\frac{4-k^2}{3k^2}}$ and
    $\sin(\alpha)=\sqrt{\frac{4-k^2}{3}}$.  \hint{Solving for one
      variable in one equation and substituting into the other is
      very difficult. Instead try squaring both equations and add
      them together? From there solving for $\sin(\beta)$ is
      straightforward\footnote{Remember that ``straightforward''
        does not mean easy.}.}
  \item Use the results from part (b) to fill in the following
    table.
    \begin{table}[h]
      \centering
      \begin{tabular}{|c|c|c|c|c|c|c|} \hline
        Color& k& $\beta$ in degrees&$\alpha$ in degrees&$D(\alpha)$ in degrees&$R(\alpha)$ in degrees\\\hline
        Red & 1.331 & &&&\\\hline
        Orange & 1.332 & &&&\\\hline
        Yellow & 1.333 & &&&\\\hline
        Green & 1.335 & &&&\\\hline
        Blue & 1.337 & &&&\\\hline
        Indigo & 1.340 & &&&\\\hline
        Violet & 1.344 & &&&\\\hline
      \end{tabular}
    \end{table}
  \end{enumerate}
\end{embeddedproblem}

% \begin{wrapfigure}[]{r}{1.25in}
%   \captionsetup{labelformat=empty}
%   \centerline{\includegraphics*[height=.in,width=1.25in]{../Figures/DoubleRainbow7}}
%   \label{fig:DoubleRainbow7}
% \end{wrapfigure}
The values of $R(\alpha)$ you obtained in the table
accounts for the various color bands that appear in the rainbow as
sunlight hits droplets at various angles.% \\
% \centerline{\includegraphics*[height=1.in,width=3in]{../Figures/DoubleRainbow7}}

The primary rainbow is created by light entering the upper portion of
the raindrop.  If light is bright enough then a secondary rainbow is
created from light entering the lower portion of the raindrop as in
the diagram below.\\
\centerline{\includegraphics*[height=2in,width=3.3in]{../Figures/DoubleRainbow8}}
This time the deviation angle $D(\alpha)$ is the total amount of counterclockwise
rotation and the rainbow angle would be $R(\alpha)=D(\alpha)-\pi$.  Note that the extra
reflection accounts for the rainbow being not as bright.

\begin{embeddedproblem}
  Show that in the case of the secondary rainbow
  $$
  D(\alpha)=2\pi+2\alpha-6\beta
  $$
  and show that this is minimized when
  $\sin(\beta)=\sqrt{\frac{9-k^2}{8k^2}}$ and
  $\sin(\alpha)=\sqrt{\frac{9-k^2}{8}}$.
  Use this to fill in the following table:
  \begin{table}[h]
    \centering
    \begin{tabular}{|c|c|c|c|c|c|c|} \hline
      Color& k& $\beta$ in degrees&$\alpha$ in degrees&$D(\alpha)$ in degrees&$R(\alpha)$ in degrees\\\hline
      Red & 1.331 & &&&\\\hline
      Orange & 1.332 & &&&\\\hline
      Yellow & 1.333 & &&&\\\hline
      Green & 1.335 & &&&\\\hline
      Blue & 1.337 & &&&\\\hline
      Indigo & 1.340 & &&&\\\hline
      Violet & 1.344 & &&&\\\hline
    \end{tabular}
  \end{table}
\end{embeddedproblem}

A couple of comments about the rainbow problem are in order.

First, notice that although Calculus was crucial to the solution it wasn't
the \emph{whole} solution. Indeed, when we look back it seems that we
hardly used Calculus at all. We found the derivatives of $D(\alpha)$
and  $R(\alpha)$ and we used Newton's Method to find the roots of those
derivatives. And that's it. That's all of the Calculus we used.

Calculus was crucial, to be sure, but it is only one tool in our
toolbox. This will generally be true of the problems you encounter in
the future. There is a tendency to want to do ``Calculus'' because
this is a ``Calculus'' course. But this is a \underline{Mathematics}
course, first and foremost. You will need all of the mathematical
tools you have. 

% Second, in our comments on Proficiency~Drill~\ref{PD:DerivUnd} we
% observed that  Fermat's Theorem specifically requires that a function
% be differentiable at the point in question. This may seem like a
% pointless formality. Obviously, if the derivative doesn't exist at
% $x=a$ then we can't conclude that it is equal to zero at $x=a$. It
% hardly seems worth the effort of writing it down.

% However, this question of differentiability is more important than it
% appears to be right now. The fact is that not all functions have
% derivatives, and some functions have tangent lines but are not
% differentiable. We have glossed over this truth so far because we were
% learning the basics, but we are about to move beyond the basics so it is
% time to begin taking explicit notice of when a derivative exists, and
% when it does not. As a first approximation we will use this
% definition:

% \begin{mydefinition}{Non-Differentiability, First Approximation}
%   The derivative of $y=f(x)$ does not exist for any value of $x$ for
%   which $\dfdx{y}{x}$ is undefined.
%   \comment{As a practical matter this means
%   any value of $x$ that would result in dividing by zero.}
% \end{mydefinition}

% As we said this is a first approximation. This is a more subtle
% problem than it appears to be at first. We will refine this definition
% later.
\begin{embeddedproblem}
  Solve the following optimization problems. (That is find the maximum
  and minimum of the objective function, subject to the constraint.)
  Without context, you need to be mindful of where $\dfdx{f}{x}$ does
  and does not exist as a location for a possible extremum.
  \begin{enumerate}[label={\bf{}(\alph*)}]
    \begin{multicols}{2}
    \item Objective function: $f=2x+y $\\
      Constraint: $x^2+2y^2=2 $
    \item Objective function: $f=xe^y $\\
      Constraint: $x^2+y^2=2 $
    \item Objective function: $f=x+y $\\
      Constraint: $xy=4 $
    \item Objective function: $f=xy $\\
      Constraint: $x^2+9y^2=18 $
    \end{multicols}
  \end{enumerate}
\end{embeddedproblem}

% Ultimately, we will need conditions under which we can
% guarantee that the derivative exists, but for now we will simply take
% explicit notice  that sometimes it doesn't by stating differentiabilty
% as a requirement, when it is needed.



\begin{wrapfigure}[]{r}{2in}
\captionsetup{labelformat=empty}
\centerline{\includegraphics*[height=.9in,width=2in]{../Figures/Refraction11}}
\label{fig:Refraction11}
\end{wrapfigure}
Finally, when we graphed $F(\alpha)$ above we only graphed the part
where $0\le\alpha\le\frac{\pi}{2}$ because those are the only angles
that make sense in this problem. However if we were to graph it for a
wider range of values of $\alpha$ we would have a graph like the one
at the right. Notice that this graph has a maximum as well as a
minimum, but our procedure only tells us if we have a horizontal
tangent line. It does not distinguish between a horizontal tangent at a
minimum from a horizontal tangent at a maximum. If we had not known,
\emph{a priori}, that the minimum was between $0$ and $\frac{\pi}{2}$
how would we have known if we'd somehow found the maximum rather than
the minimum.

This seems like a minor problem at first, but it can become quite
problematic in more complex, or more abstract problems where we do not
have as much physical or geometric intuition to guide us as we had in
the rainbow problem. We will need some way to distinguish between a
maximum and a minimum analytically, without relying on pictures or our
intuition.


% \includegraphics*[height=3in,width=5in]{../Figures/DoubleRainbow1}

% \includegraphics*[height=3in,width=5in]{../Figures/DoubleRainbow0}



%%% Local Variables: 
%%% mode: latex
%%% outline-minor-mode: t
%%% eval:(flyspell-mode)
%%% TeX-master: "DiffCalc"
%%% End: 

    \subsection{Rainbows}\label{subsec:rainbows}

    Did you ever notice a double rainbow? \\

    \begin{figure}[h]
    \captionsetup{labelformat=empty}
    \label{fig:DoubleRainbow9}
    \centerline{\includegraphics*[height=3in,width=5.6in]{../Figures/SavageEnariRainbow}}
    \caption{\textcolor{blue}{The image  of a double rainbow
        over Niagra Falls (left) is courtesy of Mary
        Savage. The image over Tsfat/Safed (right), Israel is courtesy of
        Elizabeth Eschel.}}
  \end{figure}


  Look carefully at the double rainbow in the pictures above\footnote{https://mrspa.files.wordpress.com/2010/10/i-walk-17.jpg}.
  Do you notice the the order of the colors is reversed from the
  primary (brighter) to the secondary bow?  Why do you think this
  happens?  Other questions come readily to mind.  Why is the
  primary rainbow below the secondary one? Why is the primary rainbow
  brighter and why isn't there always a secondary rainbow? 

  To tackle these questions, the first thing we need to do is to
  transform these non-mathematical questions into a mathematical
  problem.  That is, we need to build a mathematical model of the
  problem. As you will see, developing the model is probably the most
  challenging part.  Once we’ve done this, you will see that the
  Calculus part makes up a relatively small, but crucial part, of the
  solution.  But first we need to apply some geometry and use what we
  already know about refraction and reflection.

  When a ray of light enters a raindrop (which is spherical), it gets
  both refracted and reflected.\\
  \begin{figure}[h]
\centerline{\includegraphics*[height=2.4in,width=4.2in]{../Figures/Refraction7}}
  \caption{\textcolor{blue}{}}
  \label{fig:DoublePath}
\end{figure}

As you can see in the diagram, the ray of light entering the raindrop
at an angle of radian measure $\alpha$ is refracted to an angle of radian
measure $\beta$.  This is governed by Snell's Law
$\frac{\sin(\alpha)}{v_1}=\frac{\sin(\beta)}{v_2}$, where $v_1$  and $v_2$ are the
velocities of light in air and water, respectively.  Actually, only a
portion is refracted as the other portion is reflected off the
exterior of the raindrop.  Once inside the raindrop, it reflects off
of the back of the raindrop (again only a portion) so that the angle
of incidence equals the angle of reflection.  A portion of the light
then exits the raindrop governed by Snell’s Law again.  The total
amount of (clockwise) rotation in this process is called the deviation
angle $D(\alpha)$.  Notice that this total rotation is a sum of three separate
rotations.

\begin{figure}[h]
\centerline{\includegraphics*[height=2.4in,width=4in]{../Figures/Droplet1}}
    \caption{\textcolor{blue}{}}
    \label{fig:Droplet1}
  \end{figure}
  
% As you can see in the diagram in figure~\ref{fig:DoublePath}, when a  ray
% of light enters raindrop at point $A$ at an angle of radian measure $\alpha$ two
% things happen. Some of the light passes into the interior of the
% raindrop and 
% is refracted to an angle of radian measure $\beta$.  This is governed
% by Snell's Law. The rest of the light is reflected off of the raindrop
% at the same angle, $\alpha$. The reflected portion plays no part in
% our analysis.

% Once inside the raindrop, the refracted portion then strikes the back
% of the drop at point $B$ where it is split again.  The reflected
% portion then proceeds forward to point $C$ where it splits once
% more. The refracted portion then combines with light from millions of
% other droplets to form the rainbow.

% The total amount of (clockwise) rotation in this process is called the
% \term{deviation angle} $D(\alpha),$ and depends only on the angle of entry,
% $\alpha$.  Notice that this total rotation is a sum of the three separate
% rotations shown in figure~\ref{fig:Droplet1} at $A$, $B$, and $C$ as
% showin in Figure~\ref{fig:Droplet1}.

\begin{embeddedproblem}
  Show that $D(\alpha)=\pi-4\beta+2\alpha$.
\end{embeddedproblem}

Remember that $\alpha$ and $\beta$ are not independent. They are
related by Snell's Law which we will rewrite as
$$
\sin(\alpha)=\frac{v_1}{v_2}\sin(\beta)=k\sin(\beta),
$$
where $k=\frac{v_1}{v_2}$ is called \term{index of refraction} (for
water). In the case of white light, this index of refraction is about 
$\frac43$. Graphing $D(\alpha)$ for this value of $k$ and
$0\le\alpha\le\frac{\pi}{2}$ we have 
%\begin{wrapfigure}[]{l}{2in}
%  \captionsetup{labelformat=empty}
  \centerline{\includegraphics*[height=2in,width=5in]{../Figures/DoubleRainbow3}}

%   \caption{\textcolor{blue}{}}
%   \label{fig:DoubleRainbow3}
% \end{wrapfigure}
Notice from the first graph that the minimum of this curve appears to
be around the point $(1.03,2.41)$.  This says that the minimum value
of $D(\alpha)$ is about $2.41 \text{ radians }\approx138^\circ$
and it occurs when $\alpha\approx1.03 \text{ radians }\approx59.4^\circ$.  Notice from the enlargement the graph near that
minimum that for values of $\alpha$ near $1.03$ the amount of deviation is roughly
the same.  Mathematically, this was straightforward from the graph.
The real question is, ``Why is this pertinent?''  The answer is that
light rays entering the raindrop at approximately $59^\circ$ have the highest
concentration when exiting the raindrop.  At other values of $\alpha$ the
dispersion of light is greater.  It is in these concentrated light rays
that we see in the rainbow. This provides the angle that the
rainbow makes with an observer.  This supplementary angle to
$D(\alpha)$ is
called the rainbow angle,  
$$
R(\alpha)=\pi-D(\alpha)=\pi-(\pi-4\beta+2\alpha\text{ radians.}
$$
For white light we have a rainbow angle of
$$
R(1.03)=\pi-2.41\approx0.732 \text{ radians }\approx42^\circ.
$$
\centerline{\includegraphics*[height=2.5in,width=5in]{../Figures/Rainbow4}}

\begin{wrapfigure}[]{r}{2in}
\captionsetup{labelformat=empty}
\centerline{\includegraphics*[height=1in,width=2in]{../Figures/DoubleRainbow6}}
\label{fig:}
\end{wrapfigure}
To obtain the spread of colors, note that the index of refraction
changes with the color of light.  This can be seen when light shines
through a prism.

This leads to the transformation of the problem of understanding
rainbows to the mathematical problem of minimizing the objective
function subject to the constraint

\begin{embeddedproblem}
  Differentiate the objective function and the constraint to show that
  $D$ is minimized when $\cos(\alpha)=\frac{k}{2}\cos(\beta).$
\end{embeddedproblem}


When we combine this equation with our constraint, we find that we
need to solve the following system of two equations in two unknowns,
\begin{align*}
  \cos(\alpha)&=\frac{k}{2}\cos(\beta)\\
  \sin(\alpha)&=k\sin(\beta).
\end{align*}

Normally to solve such a system, we would use one of the equations to
eliminate one of the variables in the other equation.  This could be
done here there is a trick that will be cleaner.  Suppose we square
the quantities in each equation and add them together. This gives

$$
1=\cos^2(\alpha)+\sin^2(\beta)=\frac{k^2}{4}\cos^2(\beta)+k^2\sin^2(\beta).
$$
You might want to file this trick in your memory.  It comes in handy
when dealing with trig to look for sums of the squares of sines and
cosines.


\begin{embeddedproblem-enumerate}
  \begin{enumerate}[label={\bf{}(\alph*)}]
    \item Use the above equation to show that for an index of
      refraction $k$ that $D(\alpha)$  is minimized when
      $\sin(\beta)=\sqrt{\frac{4-k^2}{3k^2}}$ and
      $\sin(\alpha)=\sqrt{\frac{4-k^2}{3}}$.
  \item Use this to fill in the following table:
        \begin{table}[h]
          \begin{center}
            \begin{tabular}{|c|c|c|c|c|c|c|} \hline Color& k& $\beta$
              in degrees&$\alpha$ in degrees&$D(\alpha)$ in
              degrees&$R(\alpha)$ in degrees\\\hline Red & 1.331 &
              &&&\\\hline Orange & 1.332 & &&&\\\hline Yellow & 1.333
              & &&&\\\hline Green & 1.335 & &&&\\\hline Blue & 1.337 &
              &&&\\\hline Indigo & 1.340 & &&&\\\hline Violet & 1.344
              & &&&\\\hline
            \end{tabular}
          \end{center}
    \item  How does  this  explain the order of the colors in the
      primary rainbow?
    \end{table}

  \end{enumerate}
\end{embeddedproblem-enumerate}
% \begin{wrapfigure}[]{r}{2in}
% \captionsetup{labelformat=empty}

% \label{fig:}
% \evalues of $R(\alpha)$ you obtained in the table accounts for the various
color bands that appear in the rainbow as sunlight hits droplets at
various angles.\\
\centerline{\includegraphics*[height=1.5in,width=3in]{../Figures/Rainbow10}}

The primary rainbow is created by light entering the upper portion of
the raindrop.  If light is bright enough then a secondary rainbow is
created from light entering the lower portion of the raindrop as in
the diagram below.\\
\centerline{\includegraphics*[height=2in,width=2.6in]{../Figures/Droplet2}}
This time the deviation angle $D(\alpha)$ is the total amount of counterclockwise
rotation and the rainbow angle would be $R(\alpha)=D(\alpha)-\pi$.  Note that the extra
reflection explains why the secondary rainbow is not as bright.

\begin{embeddedproblem-enumerate}
  \begin{enumerate}[label={\bf{}(\alph*)}]
    \item Show that in the case of the secondary rainbow,
      $D(\alpha)=2\pi+2\alpha-6\beta$.
    \item Show that $D(\alpha)$ is minimized when
      $\sin(\beta)=\sqrt{\frac{9-k^2}{8k^2}}$ and
      $\sin(\alpha)=\sqrt{\frac{9-k^2}{8}}$.
    \item Use this to fill in the following table:
      \begin{table}[h]
        \begin{center}
          \begin{tabular}{|c|c|c|c|c|c|c|} \hline Color& k& $\beta$ in
            degrees&$\alpha$ in degrees&$D(\alpha)$ in
            degrees&$R(\alpha)$ in degrees\\\hline Red & 1.331 &
            &&&\\\hline Orange & 1.332 & &&&\\\hline Yellow & 1.333 &
            &&&\\\hline Green & 1.335 & &&&\\\hline Blue & 1.337 &
            &&&\\\hline Indigo & 1.340 & &&&\\\hline Violet & 1.344 &
            &&&\\\hline
          \end{tabular}
        \end{center}
      \end{table}
    \item Use this table and the table in the previous problem to
      explain why the secondary rainbow is above the primary rainbow
      and its colors are reversed.  
  \end{enumerate}
\end{embeddedproblem-enumerate}


% %%%%%%%%%%%%%%%%%%%%%%%%%%%%%%%%%%%%%%%%%%%%%%%%%%%%%%%%%%%%%%%%%%%%%%%%%
% \begin{ProficiencyDrill}
%   Use Newton's Method to show that, to two decimals, the coordinates
%   of the minimum point on the graph of $D(\alpha)$ 
%   are $(1.03, 2.41)$.

%   \hint{First you'll have to get the deviation angle in terms of
%     $\alpha$ alone.}
% \end{ProficiencyDrill}


% % \begin{wrapfigure}[]{r}{2in}
% % \centerline{\includegraphics*[height=1in,width=2in]{../Figures/Rainbow4}}
% % \caption{\textcolor{blue}{}}
% % \label{fig:Rainbow4}
% % \end{wrapfigure}
% This says that the minimum value of $D(\alpha)$ is about $2.41$
% radians $\approx 138^\circ$ and it occurs when
% $\alpha\approx1.03$ radians $\approx59.4^\circ$.  Notice from
% an enlargement the graph near that minimum that for values of $\alpha$
% near $1.03$ radians the amount of deviation is roughly the
% same.

% This says that light rays entering the raindrop at approximately
% $59^\circ$ have the highest concentration when exiting the raindrop.
% At other values of $\alpha$ the dispersion of light is greater.  This
% provides the angle that the rainbow makes with an observer.  This
% supplementary angle to $D(\alpha)$ is called the rainbow angle,
% $$
% R(\alpha)=\pi-D(\alpha)=\pi-(\pi-4\beta+2\alpha)=4\beta-2\alpha\text{
%   radians}.
% $$
% For white light, we have a rainbow angle of
% $$
% R(1.03=\pi-2.41\approx.732\text{ radians}\approx42^\circ.
% $$
% % This says that the minimum value of $D(\alpha)$ is about $2.41$
% % radians or approximately $138^\circ$  and it occurs when
% % $\alpha\approx1.03$ radians, or about $59.4^\circ$. 

% % Now notice the rectangle drawn around the minimum point in
% % Figure~\ref{fig:DoubleRainbow3}. Because it is wider than it is tall
% % we see that the broad range of incoming angles in the domain will all
% % emerge from the droplet at about the the same deviation angle,
% % approximately $138^\circ$ as seen in Figure~\ref{fig:Rainbow4}.
% Thus, where the angle between the incoming light rays and line from
% the rainbow to the observer is approximately $42^\circ$ there will be
% a bright, circular ring in the sky. This bright ring is the rainbow.

% To obtain the spread of colors, note that we said earlier that
% $k=\frac43$ was an average over all the colors in white sunlight. The
% fact that each color has its own index of refraction was a young Isaac
% Newton's for major scientific publication, \pubtitle{Optics}. This
% is the work that brought him to prominence in scientific circles.

% Newton determined experimentally that white light was a mixture of
% other colors. It was well known that the time that a beam of light
% passing through a prism would emerge with the colors separated. A prevailing
% theory of the time was that the prism had altered the light in some
% way that had added the colors. Thus it was believed that the color
% bands that emerged from the prism were a property of the prism, not of
% the light.

% Newton debunked that theory by passing \emph{only} the red band
% through a second prism. He reasoned that if the first prism had
% somehow altered the light then the second would alter it again. But
% when only red light emerged from the second prism, refracted but
% otherwise unchanged, he concluded that the color was an inherent
% property of the light, not the prism. Moreover it is because each  color
% has its own index of refraction that the colors spread out upon being
% refracted.

% To find the rainbow angle for each color we first need to solve for
% $\alpha$ and $\beta$ in terms of the refraction index, $k$.


% Normally to solve such a system, we would use one of the equations to
% eliminate one of the variables in the other equation.  This could be
% done here, but it gets very messy very fast. Fortunately there is a
% trick that will be cleaner.  Suppose we square the quantities in each
% equation and add them together. This gives
% $$
% 1=\cos^2(\alpha)+\sin^2(\beta)=\frac{k^2}{4}\cos^2(\beta)+k^2\sin^2(\beta).
% $$
% This is a very handy trick to keep in mind when dealing with
% trigonometric functions. You should file it away in your memory.

% The primary rainbow is created by light entering the upper portion of
% the raindrop.  If light is bright enough then a secondary rainbow is
% created from light entering the lower portion of the raindrop as in
% the diagram below at the right.




% % where  is the index of refraction (for water).  In the case of white light, this index of refraction is about .  Graphing  for this value of  and , we have.


% %   We will use what we
% %   know about the refraction and reflection of light to examine some of
% %   these questions.
  
% %   When light strikes the surface of water some of it reflects -- that is
% %   why you can see yourself in smooth surface of a puddle -- and some
% %   of it passes into the water, refracting as it does. It is the
% %   combination of the reflection and refraction of sunlight striking
% %   the water droplets left in the atmosphere after a rainstorm that
% %   causes rainbows.

% %   Since the surface of a water droplet is curved\aside{We will assume
% %     all water droplets are spherical.} rather than flat it
% %   would appear that the first question to address is,
% %   ``Does light  reflect or refract differently when it strikes a
% %   curved surface than when it strikes a flat one.'' The Principle of
% %   Local Linearity resolves this for us completely. Since a given beam
% %   of the light strikes the surface of the water droplet at a single
% %   point (essentially) it will reflect or refract exactly as if it had
% %   struck a plane tangent to the droplet at that point.

% %   or issue we need to  However we will find it useful to discuss the
% % reflection and refraction of light when it strikes a curved mirror.

% % \begin{wrapfigure}[]{r}{2in}
% %   \captionsetup{labelformat=empty}
% %   \centerline{\includegraphics*[height=1in,width=2in]{../Figures/Refraction10}}
% %   \label{fig:Refraction10}
% % \end{wrapfigure}
% % Happily nothing changes much. We need only observe that light only
% % interacts with the mirror locally, where it actually strikes  the
% % mirror. And since we know that at very small scales any curve is
% % indistinguishable from it's tangent line it should be clear that when
% % a light beam strikes a curved mirror it reflects or refracts exactly
% % as if it were striking a flat mirror tangent to the curve at the point
% % of interaction as seen at the right. This is again the principle of Local
% % Linearity\notetoself{We should state this as a principle at the very
% %   beginning of chapter~\ref{chapt:differentials} (or possibly before)
% %   and refer to it formally whenever possible.  -- Sep  6, 2020}  

% % \begin{wrapfigure}[]{l}{2.1in}
% %   \captionsetup{labelformat=empty}
% %   \centerline{\includegraphics*[height=.75in,width=2.1in]{../Figures/DoubleRainbow6}}
% %   \label{fig:DoubleRainbow6}
% % \end{wrapfigure}
% % The ordering of the colors of a rainbow, the shape of a
% % rainbow, its position relative to the viewer, and the simple fact that
% % rainbows exist to be seen are a result of the reflective and
% % refractive properties of light. Even better for our purposes
% % optimization is needed to answer them.


% % % wrote a seminal work on optics and .  This can be seen as light passes
% % % through a prism or in a rainbow.  Surely you heard the mnemonic
% % % \alert{ROY G BIV} for the colors of a rainbow.  Which color is on top,
% % % red or violet?  Before you go further, take a guess.  % \\
% % Look closely at the three rainbows in the images at the right above. They
% % share certain properties that we would like to understand. These are:
% % \begin{enumerate}
% % \item They are bright. At least they are brighter than the surrounding
% %   viewscape. Indeed it is only the brightness of the light coming from
% %   the apparent location of the rainbow that makes it visible since a
% %   rainbow is not a physical object.
% % \item They have a particular shape, but what is it? Are they
% %   circuluar? Parabolic? Elliptical? Some other, possibly unnamed
% %   shape?
% % \item Two of the rainbows shown are double rainbows. Notice that the
% %   secondary rainbow is dimmer\footnote{That's actually what makes it
% %     the secondary.}. Also the colors are reversed from the primary to
% %   the secondary. We'd like to understand that.
% % \end{enumerate}
% % Each of the properties listed above can be couched as an optimization
% % problem. We will address them one at a time.

% % To keep things as simple as possible we will assume that water
% % droplets suspended in the atmosphere assume a perfectly spherical
% % shape. This is close enough to reality that the errors in our
% % computations will not be significant.


% % \begin{wrapfigure}[]{l}{1.2in}
% %   \captionsetup{labelformat=empty}
% %   \centerline{\includegraphics*[height=1.2in,width=1.2in]{../Figures/Refraction6}}
% %   \label{fig:Refraction6}
% % \end{wrapfigure}
% % Also, we will analyze this as if everything were two-dimensional. That
% % is we will draw the light beam crossing into a circular raindrop
% % rather than a spherical raindrop. Again this is just to keep things
% % simple. The analysis of the three dimensional problem  yields
% % similar results but is much harder algebraically.

% % As we've seen when light refracts across the plane dividing air and
% % water it bends toward the normal (perpendicular). The physical reason
% % for this is that the velocity of light in air is faster than its
% % velocity in water. If it transitions from a slower medium (like water)
% % to a faster medium (like air) then it bends away from the normal. Thus
% % in the diagram at the right we can deduce that the velocity is lower
% % in the region below the boundary (in black) whichever direction the
% % light ray (in red) is moving.



% % % Of course the sketch above doesn't capture the curvature of the rain
% % % droplet.  Our sketch should look more like the figure at the left,
% % % where the light beam is crossing the black, circular boundary of the
% % % droplet. Does Snell's Law still apply when the boundary is curved?

% % % Of course it does. The light beam doesn't care what shape the boundary
% % % has away from the point of entry. It only interacts with the boundary
% % % locally, and locally any curve is indistiquishable from it's tangent
% % % line so the light beam will refract as it crosses into the droplet
% % % just as if it were crossing the line tangent to the circle at the
% % % point of intersection. As we saw in\notetoself{We need to add this
% % % problem somewhere before this, if we don't already have it: Show that the line tangent to a
% % % circle is orthogonal to the radius of the circle.  -- Sep 3, 2020}
% % % problem~\ref{} the line tangent to a circle at a point is orthogonal
% % % to the radius that passes through the same point. The radius and
% % % tangent are the blue, dotted lines in the diagram at the left.
% % But the physics involved is not our concern. All we need is  the
% % knowledge that Snell's Law is in play in the diagram at the left.
% % $\frac{\sin(\alpha)}{v_1}=\frac{\sin(\beta)}{v_2}.$
% % The green line in the diagram is the path the light beam would have
% % followed had it \emph{not} refracted at the boundary. 
% % \begin{ProficiencyDrill}
% %   For later
% %   reference notice that the angle between the red and green lines is
% %   $\alpha-\beta$ as indicated. \hint{This is an elementary  result
% %     from geometry. Don't  to make hard.}
% % \end{ProficiencyDrill}
% % We now have everything we need to trace the path of a beam of light
% % into and back out of a water droplet.
% % In the sketch below the red arrows indicate  the path the light
% % beam actually follows. The green arrows indicate the path that the
% % light beam would follow if it were not refracted or reflected, $O$ is
% % the center of the droplet, and the line segments $\overline{OA},
% % \overline{OB},$ and $\overline{OC}$ are all radii of the droplet.

% % % \begin{figure}[h]\centering
  
% % %   \centerline{\includegraphics*[height=2.8in,width=5in]{../Figures/Refraction7}}
  
% % %   \caption{The path of a typical light beam as it moves through a
% % %     spherial droplet.}
% % %   \label{fig:DoublePath}
% % % \end{figure}

% % In order to address the question of the rainbow's brightness -- why it
% % is even visible in the first place -- we want to find the total
% % clockwise rotation of the direction of the beam from the time it
% % enters the raindrop to the time it exits. That is, if we place the
% % starting and ending beams head-to-tail as indicated in the sketch at the right we
% % want to find the angle $D(\alpha)$.

% % \begin{wrapfigure}[]{r}{1.5in}
% %   \captionsetup{labelformat=empty}
% %   \centerline{\includegraphics*[height=.6in,width=1.5in]{../Figures/Refraction8}}
% %   \label{fig:Refraction8}
% % \end{wrapfigure}
% % This is called the
% % \term{deviation angle} and for a given entry angle $\alpha$ it tells
% % us what the exit angle will be. We are interested in finding the value
% % of $\alpha$ that \emph{minimizes} $D(\alpha)$.

% % Here's why. Suppose $D(\alpha)$ is mimimum for some angle, say
% % $\alpha=\phi$, and that $D(\phi)=\psi$ is the minimum. Then all of the
% % light beams entering the droplet at an angle \emph{near} $\phi$ will
% % exit \emph{near} $\psi$.

% % This is always true of course. Things that start out together
% % generally end up together. But the nature of the function
% % $D(\alpha)$ guarantees that there is a large interval of entry
% % angles which will all have approximately the same exit angle.  The
% % easiest way to see this is to look at the the shape of the graph of
% % $D(\alpha)$ below. Notice that the
% % black box we drew around then minimum point (denoted\footnote{For now you
% %   will have to trust that we have represented $D(\alpha)$
% %   accurately. We will shortly be deriving a formula for $D(\alpha)$
% %   and you can verify the accuracy of this graph yourself.}
% % $(\gamma,\psi)$ in the sketch.). is much wider than it is
% % tall.
% % % This says that all of the the light beams entering at an angle
% % % between about $0.8$ radians and $1.2$ radians will exit at about $2.4$
% % % radians. That is the light will be concentrated (brighter) at a
% % % deviation angle of about $2.4$ radians (about $138^\circ$. 

% % % We will soon see that the deviation angle
% % % depends \emph{only} on the angle $\alpha$ as we've
% % % indicated. 
% % In fact from the dotted lines and the scales on the axes we can easily
% % see that for entry angles betweeen a little more than $0.8\approx45.8^\circ$ radians and
% % a little more than $1.2\approx68.8^\circ$ radians the exit angles will all be very
% % close to $2.4\approx138^\circ$ radians. Outside of the interval $(.8, 1.2)$ the
% % exit angles grow very quickly.

% % This says that the light beams coming out of a droplet at an angle of
% % $\psi\approx2.4$ radians are a concentration of the beams that entered
% % between about $.8$ and $1.2$ radians. So light
% % exiting at about $2.4$ radians will be brighter than the surrounding
% % light. This is what makes a rainbow visible in the first place.

% % Now of course, we need to do the calculations to confirm all of
% % this. First we need a formula for the deviation angle, $D(\alpha)$
% % so we will follow the path of a typical light beam as it refracts and
% % reflects its way into and then back out of  a typical droplet.  

% % In Figure~\ref{fig:DoublePath} we see that the beam (in red) makes a
% % turn of $\alpha-\beta$ at the point $A$ and another turn of the
% % same magnitude at point $C$ so
% % $$
% % D(\alpha) =
% % 2(\alpha-\beta) + \{\text{the magnitude of the rotation at point
% %   $B$}\}.
% % $$

% % \begin{wrapfigure}[]{r}{2in}
% %   \captionsetup{labelformat=empty}
% %   \centerline{\includegraphics*[height=1.8in,width=2in]{../Figures/Refraction9}}
% %   \label{fig:Refraction9}
% % \end{wrapfigure}
% % To find the magnitude of the rotation at point $B$ consider the sketch at
% % the right. Clearly at point $B$ the beam turns by the angle labeled
% % $\theta$.

% % \begin{ProficiencyDrill}
% %   \begin{enumerate}[label={\bf{}(\alph*)}]
% %   \item   Use Geometry to show that $\alpha=\angle AOB$.
% %   \item Show that $\alpha=\pi-2\beta$. \hint{The interior angles
% %       of triangle $OAB$ must sum to $\pi$. What is $\beta+\theta$
% %       equal to?}
% %   \end{enumerate}
% % \end{ProficiencyDrill}
% % Therefore
% % \begin{align*}
% %   D(\alpha) &= 2(\alpha-\beta) + \{\text{the magnitude of the
% %               rotation at point  $B$}\}\\
% %             &= 2(\alpha-\beta) + \alpha\\
% %             &= 2(\alpha-\beta) + \pi-2\beta, \text{ or}\\
% %             &= \pi+2\alpha-4\beta.
% % \end{align*}

% % Notice that $\alpha$ and $\beta$ are not independent, they are governed by Snell's
% % Law which we will rewrite as
% % \begin{equation}
% %   \sin(\alpha)=\frac{v_{\alpha}}{v_{\beta}}\sin(\beta)=k\sin(\beta).\label{eq:IndRef}
% % \end{equation}
% % The constant $k=\frac{v_{\alpha}}{v_{\beta}}$ is called  the ``index of refraction'' (for
% % water).
% % \begin{ProficiencyDrill}
% %   Show that $\beta=\inverse\sin\left(\frac{\sin(\alpha)}{k}\right)$.
% % \end{ProficiencyDrill}
% % Knowing $\beta$ we can now express the deviation angle,
% % $D(\alpha)$ as a function of $\alpha$ alone:
% % \begin{align*}
% %   D(\alpha) &= \pi+2\alpha-4\beta \\
% %   D(\alpha) &= \pi+2\alpha-4\inverse\sin\left(\frac{\sin(\alpha)}{k}\right).
% % \end{align*}

% % \begin{embeddedproblem}
% %   For visible light  the index of refraction
% %   is $k\approx\frac43$. Mathematically, the problem becomes minimizing
% %   $D(\alpha)=\pi+2\alpha-4\beta$ subject to the constraint
% %   $\sin(\alpha)=k\sin(\beta)$, $ 0\le\alpha<\frac{\pi}{2} $.

% %   \begin{enumerate}[label={\bf{}(\alph*)}]
% %   \item Graph $D(\alpha)$ with $k=\frac43$ and confirm that the graph
% %     we showed you earlier is correct.
% %   \item Use Newton's Method to show that the coordinates,
% %     $(\gamma,\psi),$
% %     of the minimum point on the graph are $\gamma\approx1.04\approx59.59^\circ$ and
% %     $\psi\approx2.41\approx138.1^\circ$ respectively.
% %   \end{enumerate}
% % \end{embeddedproblem}
% % % ****************** Pick up here *************************************

% % % Graphing $D(\alpha)$ for this value of  $k$ and
% % % $0\le{}\alpha{}<\frac{\pi{}}{2}$, we have.
% % % Notice from the graph that the minimum of this curve appears to be
% % % around the point $(1.03, 2.41)$.  This says that the minimum value of
% % % $D(\alpha)$  is about $2.41\text{ radians }\approx 138^\circ$  and it
% % % occurs when $\alpha\approx1.03\text{ radians }\approx59.4^\circ$.
% % % As we can see from our sketch the  minimum 
% % % value of $\alpha$ near  $1.03$  the amount of deviation is roughly
% % % the same.

% % % \vskip.5in{}
% % % ***************************************************************\\
% % % As you can see in the diagram below, the ray of light entering the
% % % raindrop at an angle of radian measure $\alpha$ is refracted to an
% % % angle of radian measure $\beta$.  This is governed by Snell's Law.\\
% % % % \begin{wrapfigure}[]{l}{3in}
% % % %   \captionsetup{labelformat=empty}
% % % %   \centerline{\includegraphics*[height=1.75in,width=3.25in]{../Figures/DoubleRainbow2}}
% % % %   \label{fig:DoubleRainbow2}
% % % % \end{wrapfigure}
% % % $\frac{\sin(\alpha)}{v_\alpha}=\frac{\sin(\beta)}{v_\beta}$, where
% % % $v_\alpha$ and $v_\beta$ are the velocities of light in air and water,
% % % respectively.  Actually, only a portion is refracted as the other
% % % portion is reflected off the exterior of the raindrop.  Once inside
% % % the raindrop, it reflects off of the back of the raindrop (again only
% % % a portion) so that the angle of incidence equals the angle of
% % % reflection.  A portion of the light then exits the raindrop governed
% % % by Snell's Law again.  The total amount of (clockwise) rotation in
% % % this process is called the deviation angle, $D(\alpha)$.%

% % % \begin{embeddedproblem}
% % %   Show that $D(\alpha)=\pi-4\beta+2\alpha.$
% % % \end{embeddedproblem}


% % \begin{wrapfigure}[]{r}{2.5in}
% %   \captionsetup{labelformat=empty}
% %   \includegraphics*[height=.8in,width=2.5in]{../Figures/DoubleRainbow5}
% %   \label{fig:DoubleRainbow5}
% % \end{wrapfigure}
% % This says that light rays entering the raindrop at approximately
% % $59^\circ$ have the highest concentration when exiting the raindrop.
% % At other values of $\alpha$ the dispersion of light is greater.  This
% % provides the angle that the rainbow makes with an observer.  This
% % supplementary angle to $D(\alpha)$ is called the rainbow angle,
% % $$
% % R(\alpha) =\pi-D(\alpha)=\pi-(\pi+2\alpha-4\beta)=4\beta-2\alpha\text{
% %   radians.}
% % $$
% % For white light, we have a rainbow angle of
% % $$
% % R(1.03)\approx\pi-2.41\approx0.734\text{ radians}\approx42^\circ.
% % $$

% % To obtain the spread of colors, note that the index of refraction
% % changes with the color of light.  We said earlier that the index of
% % refraction for visible light is approximately $\frac43$. This is true,
% % but it actually varies with the frequency (color) of the light. The
% % red light refracts at one angle and the orange at a slightly different
% % angle so they appear in the rainbow adjacent to each other. This
% % variation is continuous throughout the spectrum which is what gives us
% % the rainbow.


% % \begin{embeddedproblem}
% %   We could use Newton's Method like we did earlier to find the
% %   minimum of $D(\alpha)$ for various values of $k$ but it is a
% %   little easier to do it this way:
% %      \begin{enumerate}[label={\bf{}(\alph*)}]
% %   \item Show that $D(\alpha)$ is minimized when % $\cos(\alpha)=\frac{k}{2}\cos(\beta)$.
% %   \item Use the equation obtained in part (a) along with the
% %     constraint $\sin(\alpha)=k\sin(\beta)$ to show that $D(\alpha)$
% %     is minimized when $\sin(\beta)=\sqrt{\frac{4-k^2}{3k^2}}$ and
% %     $\sin(\alpha)=\sqrt{\frac{4-k^2}{3}}$.  \hint{Solving for one
% %       variable in one equation and substituting into the other is
% %       very difficult. Instead try squaring both equations and add
% %       them together? From there solving for $\sin(\beta)$ is
% %       straightforward\footnote{Remember that ``straightforward''
% %         does not mean easy.}.}
% %   \item Use the results from part (b) to fill in the following
% %     table.
% %     \begin{table}[h]
% %       \centering
% %       \begin{tabular}{|c|c|c|c|c|c|c|} \hline
% %         Color& k& $\beta$ in degrees&$\alpha$ in degrees&$D(\alpha)$ in degrees&$R(\alpha)$ in degrees\\\hline
% %         Red & 1.331 & &&&\\\hline
% %         Orange & 1.332 & &&&\\\hline
% %         Yellow & 1.333 & &&&\\\hline
% %         Green & 1.335 & &&&\\\hline
% %         Blue & 1.337 & &&&\\\hline
% %         Indigo & 1.340 & &&&\\\hline
% %         Violet & 1.344 & &&&\\\hline
% %       \end{tabular}
% %     \end{table}
% %   \end{enumerate}
% % \end{embeddedproblem}

% % % \begin{wrapfigure}[]{r}{1.25in}
% % %   \captionsetup{labelformat=empty}
% % %   \centerline{\includegraphics*[height=.in,width=1.25in]{../Figures/DoubleRainbow7}}
% % %   \label{fig:DoubleRainbow7}
% % % \end{wrapfigure}
% % The values of $R(\alpha)$ you obtained in the table
% % accounts for the various color bands that appear in the rainbow as
% % sunlight hits droplets at various angles.% \\
% % % \centerline{\includegraphics*[height=1.in,width=3in]{../Figures/DoubleRainbow7}}

% % The primary rainbow is created by light entering the upper portion of
% % the raindrop.  If light is bright enough then a secondary rainbow is
% % created from light entering the lower portion of the raindrop as in
% % the diagram below.\\
% % \centerline{\includegraphics*[height=2in,width=3.3in]{../Figures/DoubleRainbow8}}
% % This time the deviation angle $D(\alpha)$ is the total amount of counterclockwise
% % rotation and the rainbow angle would be $R(\alpha)=D(\alpha)-\pi$.  Note that the extra
% % reflection accounts for the rainbow being not as bright.

% % \begin{embeddedproblem}
% %   Show that in the case of the secondary rainbow
% %   $$
% %   D(\alpha)=2\pi+2\alpha-6\beta
% %   $$
% %   and show that this is minimized when
% %   $\sin(\beta)=\sqrt{\frac{9-k^2}{8k^2}}$ and
% %   $\sin(\alpha)=\sqrt{\frac{9-k^2}{8}}$.
% %   Use this to fill in the following table:
% %   \begin{table}[h]
% %     \centering
% %     \begin{tabular}{|c|c|c|c|c|c|c|} \hline
% %       Color& k& $\beta$ in degrees&$\alpha$ in degrees&$D(\alpha)$ in degrees&$R(\alpha)$ in degrees\\\hline
% %       Red & 1.331 & &&&\\\hline
% %       Orange & 1.332 & &&&\\\hline
% %       Yellow & 1.333 & &&&\\\hline
% %       Green & 1.335 & &&&\\\hline
% %       Blue & 1.337 & &&&\\\hline
% %       Indigo & 1.340 & &&&\\\hline
% %       Violet & 1.344 & &&&\\\hline
% %     \end{tabular}
% %   \end{table}
% % \end{embeddedproblem}

% % A couple of comments about the rainbow problem are in order.

% % First, notice that although Calculus was crucial to the solution it wasn't
% % the \emph{whole} solution. Indeed, when we look back it seems that we
% % hardly used Calculus at all. We found the derivatives of $D(\alpha)$
% % and  $R(\alpha)$ and we used Newton's Method to find the roots of those
% % derivatives. And that's it. That's all of the Calculus we used.

% % Calculus was crucial, to be sure, but it is only one tool in our
% % toolbox. This will generally be true of the problems you encounter in
% % the future. There is a tendency to want to do ``Calculus'' because
% % this is a ``Calculus'' course. But this is a \underline{Mathematics}
% % course, first and foremost. You will need all of the mathematical
% % tools you have. 

% % % Second, in our comments on Proficiency~Drill~\ref{PD:DerivUnd} we
% % % observed that  Fermat's Theorem specifically requires that a function
% % % be differentiable at the point in question. This may seem like a
% % % pointless formality. Obviously, if the derivative doesn't exist at
% % % $x=a$ then we can't conclude that it is equal to zero at $x=a$. It
% % % hardly seems worth the effort of writing it down.

% % % However, this question of differentiability is more important than it
% % % appears to be right now. The fact is that not all functions have
% % % derivatives, and some functions have tangent lines but are not
% % % differentiable. We have glossed over this truth so far because we were
% % % learning the basics, but we are about to move beyond the basics so it is
% % % time to begin taking explicit notice of when a derivative exists, and
% % % when it does not. As a first approximation we will use this
% % % definition:

% % % \begin{mydefinition}{Non-Differentiability, First Approximation}
% % %   The derivative of $y=f(x)$ does not exist for any value of $x$ for
% % %   which $\dfdx{y}{x}$ is undefined.
% % %   \comment{As a practical matter this means
% % %   any value of $x$ that would result in dividing by zero.}
% % % \end{mydefinition}

% % % As we said this is a first approximation. This is a more subtle
% % % problem than it appears to be at first. We will refine this definition
% % % later.
% % \begin{embeddedproblem}
% %   Solve the following optimization problems. (That is find the maximum
% %   and minimum of the objective function, subject to the constraint.)
% %   Without context, you need to be mindful of where $\dfdx{f}{x}$ does
% %   and does not exist as a location for a possible extremum.
% %   \begin{enumerate}[label={\bf{}(\alph*)}]
% %     \begin{multicols}{2}
% %     \item Objective function: $f=2x+y $\\
% %       Constraint: $x^2+2y^2=2 $
% %     \item Objective function: $f=xe^y $\\
% %       Constraint: $x^2+y^2=2 $
% %     \item Objective function: $f=x+y $\\
% %       Constraint: $xy=4 $
% %     \item Objective function: $f=xy $\\
% %       Constraint: $x^2+9y^2=18 $
% %     \end{multicols}
% %   \end{enumerate}
% % \end{embeddedproblem}

% % % Ultimately, we will need conditions under which we can
% % % guarantee that the derivative exists, but for now we will simply take
% % % explicit notice  that sometimes it doesn't by stating differentiabilty
% % % as a requirement, when it is needed.



% % \begin{wrapfigure}[]{r}{2in}
% % \captionsetup{labelformat=empty}
% % \centerline{\includegraphics*[height=.9in,width=2in]{../Figures/Refraction11}}
% % \label{fig:Refraction11}
% % \end{wrapfigure}
% % Finally, when we graphed $F(\alpha)$ above we only graphed the part
% % where $0\le\alpha\le\frac{\pi}{2}$ because those are the only angles
% % that make sense in this problem. However if we were to graph it for a
% % wider range of values of $\alpha$ we would have a graph like the one
% % at the right. Notice that this graph has a maximum as well as a
% % minimum, but our procedure only tells us if we have a horizontal
% % tangent line. It does not distiguish between a horizontal tangent at a
% % minimum from a horizontal tangent at a maximum. If we had not known,
% % \emph{a priori}, that the minimum was between $0$ and $\frac{\pi}{2}$
% % how would we have known if we'd somehow found the maximum rather than
% % the minimum.

% % This seems like a minor problem at first, but it can become quite
% % problematic in more complex, or more abstract problems where we do not
% % have as much physical or geometric intuition to guide us as we had in
% % the rainbow problem. We will need some way to distinquish between a
% % maximum and a minimum analytically, without relying on pictures or our
% % intuition.


% % % \includegraphics*[height=3in,width=5in]{../Figures/DoubleRainbow1}

% % % \includegraphics*[height=3in,width=5in]{../Figures/DoubleRainbow0}


\section{More Optimal Squares}
\label{sec:more-optimal-squares}
Let's get back to our original problem of finding maximum and minimum
areas of squares with one fixed corner and a second corner on a
constraint curve.  We already did this when the constraint curve was a
line, circle, or ellipse.  Things get very interesting when we look at
some other constraint curves.

\begin{myexample}[Constructing a Square on an Hyperbola]
  \label{example:HyperbolaDist1}
  Find the area of the maximal square that can be constructed with one
  corner at the point $(2,3)$ and one corner on the hyperbola:
  $$x^2-y^2=1.$$
  % \centerline{\includegraphics*[height=2in,width=3in]{../Figures/HyperbolaDist1}}

  Earlier (Problem~\ref{problem:EllipseDist}) we were careful to point
  out the value of visualizing a problem with a sketch. To emphasize
  the importance of this we will solve this problem without
  visualization. Naturally, something will go wrong. See if you can
  spot it as we proceed.

  As before the area of an arbitrary square constructed as described is
  given by
  $$
  A=(x-2)^2+(y-3)^2.
  $$
  so
  $$
  \dx{A} = 2(x-2)\dx{x} + 2(y-3)\dx{y}.
  $$
  but this time the constraint is $x^2-y^2=1$ so we have
  $\dx{y}=\frac{x}{y}\dx{y}.$ 
  Thus
  \begin{equation}
    \dfdx{A}{x} =
    2\left(2x-2-\frac{3x}{\sqrt{x^2-1}}\right)\label{eq:HypDeriv}
  \end{equation}
  
  \begin{ProficiencyDrill}
    \begin{enumerate}[label={\bf{}(\alph*)}]
    \item Use Newton's Method with an initial guess of $r_0=3$ to show
      that we have a minimal square when $x\approx 2.623.$
    \item Find the corresponding $y$ value, and show that the area of
      the corresponding optimal square is  $A(x)\approx0.719$.
    \item Show that this square is \alert{not} maximal by computing
      the area of \emph{any} other square with one corner on the
      hyperbola.
    \end{enumerate}
  \end{ProficiencyDrill}

\begin{wrapfigure}[]{r}{2in}
\captionsetup{labelformat=empty}
\centerline{\includegraphics*[height=2in,width=2in]{../Figures/HyperbolaDist6}}
\label{fig:}
\end{wrapfigure}  
We warned you that something would go
wrong. Do you see what it was?

The problem asks us to find the maximal square, but we found the
minimal.  How did that happen?

It was entirely because we failed to visualize the problem with a
sketch, so take a look at the visualization at the right.  Can you see
that there must two optimal squares? Can you also
see that corner of the minimal square on the hyperbola will have an
$x$ coordinate somewhere between $2$ and $3$. The $x$ coordinate we
found was approximately $2.623$, so clearly we found the minimal
square, not the maximal as required. The other optimal square has a corner
on the left branch of the graph not the right, so the $x$ coordinate
will be negative.



  \begin{wrapfigure}[]{l}{1.2in}
    \captionsetup{labelformat=empty}
    \vskip-5mm{}
    \centerline{\includegraphics*[height=1in,width=1.2in]{../Figures/HyperbolaDist4}}
    \label{fig:HyperbolaDist1}
  \end{wrapfigure}
We used Newton's Method to find the roots of $\dfdx{A}{x}$ because we
know that when $\dfdx{A}{x}=0$ $A(x)$ will be optimal. But recall that in
Chapter~\ref{chapt:Approximation} we saw that it can fail in a couple
of ways if we aren't careful. In the worst case if we pick an initial 
guess that is not sufficiently close to the root we seek nothing goes wrong
with our calculations, it just returns the wrong answer. That's what
happened here. We chose to start with $r_0=3$. Had we also  drawn the
graph of $\dfdx{A}{x}$ (seen at the left) we'd have known that
starting with a positive guess would fail.

Perhaps you saw these errors as we went through the problem. If so,
good for you! But ask yourself this: If we hadn't warned you that
things were going to go wrong would you have noticed the errors?
\emph{Always} sketch the problem if you possibly can. It works.  Had
we draw this sketch \emph{before} beginning to work on the problem
we'd have know that we had the wrong answer as soon as Newton's Method
gave us back a positive value of $x$.

  \begin{ProficiencyDrill}
    Now that we see that $\dfdx{A}{x}$ has a second root it is easy to
    resolve some of these questions.
    \begin{enumerate}[label={\bf{}(\alph*)}]
    \item     Take $r_0=\frac32$, and use
      Newton's Method to show that the second root is at
      $x\approx -1.313$ 
    \item Show that $y\approx 0.851$.
    \item Show that  the area of the resulting square is
      $A(x)\approx15.59$.
    \item Is this the maximal square the problem asks for?
    \end{enumerate}
  \end{ProficiencyDrill}
\begin{wrapfigure}[]{r}{2in}
\captionsetup{labelformat=empty}
  \centerline{\includegraphics*[height=1.5in,width=2in]{../Figures/HyperbolaDist1}}
\label{fig:}
\end{wrapfigure}
In the  sketch at the right we see the hyperbola and both of the squares
we have found:

Does it appear to you that one of the sides of each square is
orthogonal (perpendicular) to the hyperbola?

We'd be willing to bet that you're expecting that the next  problem
will be  to show that the lines joining the point, $(2,3)$ to the optimal
points on the hyperbola are perpendicular (orthogonal) to it,
as we did with the line and the circle.  Well, you're right. That is
the next problem.


But before we start let's pause, step back and observe that it seems
intuitively clear that this \emph{always} true.

Optimizing the area of a square is equivalent to optimizing the length
of one side of a square, and it should be intuitively clear that the
shortest or longest line from a point to a smooth curve -- regardless
of the shape of the curve -- is the one that is perpendicular
(orthogonal) to the curve.

  If this is true then it might be worth biting the bullet and working through
  the problem in full generality -- essentially do all such problems
  at once -- so that we never need to do it for a specific curve, with
  all of the computational complications that come with that.

  [This is surprisingly straightforward.  There is a little bit of a
  difficulty where the tangent line is horizontal or vertical, but
  you'll figure it out.]

\begin{embeddedproblem}
\begin{wrapfigure}[]{l}{1in}
\captionsetup{labelformat=empty}
\centerline{\includegraphics*[height=.7in,width=1in]{../Figures/MaxOrth1}}
\label{fig:}
\end{wrapfigure}
  Consider a curve $C$ and a point $(a,b)$  not on the curve.  Show
  that if the area of a square with corners $(a,b)$ and $(x,y)$  on
  the curve is minimized or maximized at $(x_0, y_0)$, then the line
  joining $(a,b)$  and $(x_0,y_0)$  is perpendicular (orthogonal) to
  the tangent line of $C$ at $(x_0,y_0)$ .

  Try to solve this problem completely on your own, but if you get
  stuck here is an outline you can follow:
  \begin{enumerate}[label={\bf{}(\alph*)}]
    \item Find the objective function. This will be the area of the
      square on the line from $(a,b)$ to $(x,y)$.
    \item Find the slope of $C$ at the optimum point, $(x_0,y_0)$.
    \item Find the  slope of the line from $(a,b)$ to $(x_0,y_0)$.
    \item Compare the two slopes.
    \item What happens if the tangent line is horizontal? If it is
      vertical?
  \end{enumerate}
\end{embeddedproblem}

  
Returning to our problem, the square at $x\approx -1.313$ is larger than the one at
  $x\approx 2.623$, so it must be a square with maximal area, right?

  No, unfortunately that isn't true either. This is actually a second
  \emph{minimal} square. It is clear that the
  ``wings'' of the hyperbola get arbitrarily far from the point
  $(2,3)$ as we allow $x$ to move to the right or left so there
  clearly is not maximal square.

  But how can there be two \emph{minimal} squares. That seems
  impossible.

  Look closely that the following closeup of the
  region around the point $(-1.313, 0.851)$.\\
  \centerline{\includegraphics*[height=1.6in,width=2in]{../Figures/HyperbolaDist2}}
  Clearly the blue line (one edge of the blue square above) is shorter
  than either of the green lines, so the area of the square generated
  from the blue line will be smaller than the area of a square
  generated from either of the green squares.

  So there are two distinct squares to be found, but both of them are
  minimal. The black one is clearly the smallest possible square we
  can build this way.  But if the blue square is larger than the
  black  in what
  sense can the blue square be considered  minimal?
\end{myexample}


\subsection*{\underline{Global vs. Local Extrema} }
In example~\ref{example:HyperbolaDist1}{} the minimal square (the
black one) has the smallest area of \emph{all possible} squares. We'll
call it the \term{global minimum}.  The area of the blue square is
larger than the area of the black one so it is not a global minimum
but its area is smaller than the area of all of the squares
 we can construct \emph{nearby}. We'll call it
a \term{local minimum}.

It is not too difficult to think about the areas of squares. That's
why we chose these problems for our first examples. But for more
real world problems (e.g. explaining the colors of a rainbow) it can
be difficult to envision what is happening physically. We would like a
visual representation of the problem that we can visualize easily,
regardless of the actual problem.

This is what the graph of a function is for.

When we plot the
graph of the objective function in Problem~\ref{example:HyperbolaDist1},
$ A(x) = (x-2)^2+(y-3)^2$ and account for the constraint:
$y=2\sqrt{x^2-1}$ the two extremal squares appear as the two low
points on the graph: $(-1.31, 15.59)$ and $(2.62, 0.72)$. The first
coordinate gives the $x$-coordinate of the constructed corner of the
square and the second coordinate gives the area. As you can see the
square at $x=2.62$ is the global minimum; it is the lowest point on
the square.

You can also see that the square at $x=-1.31$ is a \emph{local}
minimum; it is the lowest point on the graph \emph{locally} (nearby).




% \begin{wrapfigure}[]{r}{2.2in}
% \captionsetup{labelformat=empty}
\centerline{\includegraphics*[height=2.5in,width=3.5in]{../Figures/HypObjFun1}}
% \label{fig:}
% \end{wrapfigure}

\begin{ProficiencyDrill-1line}
  Explain why the graph has two parts.
\end{ProficiencyDrill-1line}

Because they are simpler, and easier to think about the graphs of
objective functions sometimes shine a light on aspects of the problem
that aren't clear otherwise. For example, do you see that this graph
also seems to have a local maximum at the point $(-1,18)$? This implies
that there is a third square with one corner at $x=-1$ and area equal
to $18$ which is larger than all of the nearby squares.

\begin{ProficiencyDrill}
  Sketch this problem again and include this third square.
\end{ProficiencyDrill}


The existence of local extrema really complicates things
for us because it means we  have to find all of the \term{local}
 extrema and figure out which of them is  global, if any.
We're going to have a lot of trouble keeping this all straight if we
don't get organized.

Let's do that, starting with a formal definition of (local and global)
extrema.

In the sketch below $y(a)$ is the \term{global} minimum of $y$ since it is the
lowest point on the graph.\\
\centerline{\includegraphics*[height=2.5in,width=3.7in]{../Figures/LocalGlobalDef}}

How can we distinguish it from the \term{local} maximum, $y(c)?$ If we
restrict our attention to the interval between the dotted lines, then
$y(c)$ is maximum on the interval $[c-\eps, c+\eps]$.  This suggests
the following definitions.

\begin{mydefinition}[Global Minimum]
  Suppose that $g$ is a number in the domain of $y$ such
  that\aside{The phrase ``such that'' is a term of art in
    mathematics. It means, ``with the following property.''}
  $y(g)\le y(x)$ for every $x$ in the domain of $y$.  Then $y(g)$ is a
  \term{global minimum} of $y$.
\end{mydefinition}

\begin{mydefinition}[Local Minimum]
  Suppose $l$ is a number in the domain of $y$.  Then $y(l)$ is a
  \term{local minimum} if there is an open interval $I$,  containing $l$, such that
  $y(l)\le y(x)$ for every $x\in I.$ 
\end{mydefinition}

\begin{embeddedproblem-enumerate}
  \begin{enumerate}[label={\bf{}(\alph*)}]
  \item Global and local maxima can be defined similarly. Provide two
    such definitions.
  \item Suppose that in the graph depicted above, the interval $I$ is an
    open interval, that is, it does not contain its endpoints.
    \begin{enumerate}
      \item     Would there be a global minimum?  Explain. 
      \item     Would there be a global maximum?  Explain.
    \end{enumerate}
  \item Suppose that in the graph depicted above, the interval $I$ is
      a closed interval, that is, it contains its endpoints.  Would
      there be a global maximum?  Explain.
    \begin{enumerate}
      \item     Would there be a global minimum?  Explain. 
      \item     Would there be a global maximum?  Explain.
    \end{enumerate}
  \end{enumerate}
\end{embeddedproblem-enumerate}

The contrast between parts (b) and (c) of the previous problem illustrate
that there is another kind of extremum we will need to consider.  If
our interval contains an endpoint, then we could have an extemum
there.  We will come back to this possibility later, but for now,
let’s focus on optimization without consideration of endpoints in our
domain.

In the diagram above $y(c)$ is a local minimum, $y(a)$ is the global
minimum and $y(b)$  is a local maximum.

\begin{ProficiencyDrill}
  Convince yourself that if a  \term{global} maximum (minimum) exists
  we can find it by 
  by finding all of the \term{local} maxima (minima) and selecting the
  largest (smallest) of these.
\end{ProficiencyDrill}

In example~\ref{example:HyperbolaDist1} we had a situation with two
local minima, one of which was also the global minimum. When we
graphed our objective function we saw that we also have a local
maximum at $x=-1$.


\begin{embeddedproblem}[Constructing a Square on a Parabola]
  \label{problem:Parab1}
  Find the area of the minimal square that can be constructed with
  one corner at the point $(2,3)$ and one corner on the parabola
  \begin{equation}
    (x+4)^2-y=3
    \label{eq:ParabolaConstraint}.
  \end{equation}
  It is evident
  from the diagram below that there is no maximal square.\\
  \centerline{\includegraphics*[height=2in,width=3in]{../Figures/ParabDist1}}
  \begin{enumerate}[label={\bf{}(\alph*)}]
  \item
    Show that  the objective function for this problem is:
    \begin{align*}
      A(x) &= (x-2)^2+(y-3)^2\\
      \intertext{so that}
      \dx{A}&=2(x-2)\dx{x}  +2(y-3)\dx{y}.
    \end{align*}
    % Our goal is to set $\dfdx{A}{x}=0$ and solve but first we will
    % need to get everything in terms of one variable or the other.
  \item Differentiate  Equation~\ref{eq:ParabolaConstraint} to show that
    $$\dx{y} = 2(x+4)\dx{x}.$$
  \item Show that $ \dfdx{A}{x}=2\left(2(x+4)^3-11x-50\right)$.  and
    use
    % \footnote{Solving this equation presents a puzzle since it is a
      % cubic equation. You learned in Algebra to solve Quadratic
      % Equations:~$ax^2+bx+c=0$, using the Quadratic
      % Formula:~$x=\frac{-b\pm\sqrt{b^2-4ac}}{2a}$, but almost
      % certainly you were never shown the analogous ``Cubic Formula''
      % for solving Cubic Equations of the form $ax^3+bx^2+cx+d=0$:
      % \begin{align*}
      %   x&=\sqrt[3]{\left(\frac{-b^3}{27a^3}+\frac{bc}{6a^2}-\frac{d}{2a}\right)+\sqrt{\left(\frac{-b^3}{27a^3}+\frac{bc}{6a^2}-\frac{d}{2a}\right)^2+\left(\frac{c}{3a}-\frac{b^2}{9a^2}\right)^3}}\\
      %    &\quad\quad+\sqrt[3]{\left(\frac{-b^3}{27a^3}+\frac{bc}{6a^2}-\frac{d}{2a}\right)-\sqrt{\left(\frac{-b^3}{27a^3}+\frac{bc}{6a^2}-\frac{d}{2a}\right)^2+\left(\frac{c}{3a}-\frac{b^2}{9a^2}\right)^3}}\\
      %    &\quad\quad\quad-\frac{b}{3a}
      % \end{align*}
      % You can see why you were not asked to memorize this formula  in
      % Algebra class.}
    Newton's Method
    to show that $x\approx-1.419$ is an approximate solution of
    $\dfdx{A}{x}=0$.
  \item Show that $A(-1.419)\approx 12.13$ is the approximate area
    of the minimal square.
  \end{enumerate}
\end{embeddedproblem}


\section{The Abstracted Problem}
\label{subsec:optim-abstr-probl}
\markboth{\sc Optimization}{\sc Abstracted Problem}

The examples we looked at in the previous three sections illustrated several
important points. Among them,
\begin{enumerate}
\item Finding the objective function, $O(x)$, is crucial. Without it we can't even
  get started.
\item Once we've found the objective function optimizing it becomes
  the more abstract problem of finding the extreme points of the graph of
  $O(x)$. Whether we  employ Newton's dynamic view
  (derivative as velocity) or Leibniz's geometric view (derivative as
  slope) of the derivative to attack this ``abstracted'' version of
  the original optimization,  it all comes down  to finding all of
  the places where the derivative is zero. These are the locations
  where an extremum \emph{might} occur. We must then determine where it really \emph{does} occur.
\item There are two kinds of extrema: local and global. Any global
  extrema, if they exist, will be the maximum and minimum of all of
  the local extrema. Thus our first task is to identify the local
  extrema. 
\end{enumerate}

But finding the local extreme points of the graph of an arbitrary function
is more complicated than we've made it appear so far. In this section
we will address this. By the end of the section we will have two
theorems that provide a roadmap for attacking this kind of problem.

\subsection{Distinguishing Maxima from Minima}
Graph the function
$$
O(x)=x^3-3x
$$
and confirm visually that it has a local maximum at $(-1, 2)$, and a
local minimum at $(1,-2)$.

This is easy when we have the graph in front of us
but in the absence of a graph to look at how could we tell which is
the maximum and which is the minimum? We'll need an \alert{analytic}
method; one that doesn't rely on pictures.


To that end, consider the following theorem from elementary Algebra.

\begin{mytheorem-enumerate-noname}
  \label{thm:FDT-lines}
  \begin{enumerate}
  \item If the slope of a line is positive then the $y$ coordinate is
    increasing.
  \item If the slope is negative then the $y$ coordinate is
    decreasing.
  \item If the slope is zero then the $y$ coordinate is not changing.
  \end{enumerate}
\end{mytheorem-enumerate-noname}

The obvious similarity between Theorem~\ref{thm:FDT-lines} and
Theorem~\ref{thm:FDT-Newton} is not a coincidence.

Lines are easy of course, because the slope of a line is constant: If
the slope of a line is $5$ here then it is the $5$ over
there, and way back there too. But we're interested in curves, and the slope of a curve
changes from point to point. Nevertheless, since $\dfdxat{O}{x}{a}$
is the slope of a curve at the point $x=a$ it is clear that if
$\dfdxat{O}{x}{a}>0$ then the curve is increasing \emph{at the point
  $x=a.$} Similarly if $\dfdxat{O}{x}{a}<0$ then the curve is
decreasing \emph{at the point $x=a.$} The above theorem thus
generalizes as follows:

\begin{mytheorem}[The First Derivative Test (Leibnizian version)]\ \\
  \label{thm:FDT-Leibniz}
  Consider the graph of $y(x)$.
  \begin{enumerate}
  \item If $\dfdxat{y}{x}{a}>0$ for every point $x=a$ in some interval
    then the graph is increasing on that interval.
  \item If $\dfdxat{y}{x}{a}<0$ for every point $x=a$ in some interval
    then the graph is decreasing on that interval.
  \item If $\dfdxat{y}{x}{a}=0$ for every point $x=a$ in some interval
    then the graph is horizontal on that interval.
  \end{enumerate}
\end{mytheorem}
Compare Theorem~\ref{thm:FDT-Leibniz} to
Theorem~\ref{thm:FDT-Newton}. Do you see that these are equivalent;
that they are really the same theorem? We just came at the idea from two
different points of view.


 The Newtonian and the Leibnizian viewpoints are equivalent and we
don't want to be tied to either,  so we will now restate this theorem
independent of the influence of Newton or Leibniz.  However we are not
abandoning either viewpoint. Both are still valid and both are still useful, and
if one or the other seems like the natural way to think about a given
problem you should do that.

What is common to both of them is the abstract concept of the increase
or decrease of $y$ (or whatever we want to name the dependent
variable) as a function of $x$ (or whatever we want to name the
independent variable). In the Newtonian version the function gives the
position of a moving point. In the Leibnizian version a graph gives a
visual representation of how the dependent variable depends on the
independent variable.

The following  more modern version \emph{of the same theorem} is
stated in terms of functions and thus subsumes both the Newtonian and
Leibnizian versions.

\begin{mytheorem}[The First Derivative Test (Function Version)]
  \label{thm:FDT}
  Suppose 
  $$
  y=f(x).
  $$
  where $f$ is some function  of $x.$
  The following statements are true:
  \begin{enumerate}
  \item If $\dfdxat{y}{x}{a}>0,$ then $f(x)$ is increasing at
    $x=a$. 
  \item If $\dfdxat{y}{x}{a}<0$, then $f(x)$ 
    is decreasing at $x=a$.
  \item If $\dfdxat{y}{x}{a}=0$ then $f(x)$ is neither increasing
    nor decreasing at $x=a$.
  \end{enumerate}
\end{mytheorem}

As we have just seen the derivative gives us information about curves
in much the same way that the slope of a straight line gives us
information about the line. The only difference is that since the
numerical value of the derivative changes at each point it only gives
us \emph{local} information about the curve at that point. 
Moreover, since we reached the same conclusions regardless of whether
we used Newton's or Leibniz's viewpoint, we can be reasonably
confident that our conclusions are independent of which point of view we
use.

We seem to have gone pretty far afield though. We said we were going
to find an analytic way to distinguish between a maximum and a
minimum. Do you see how the First Derivative Test does that? Take a
moment to think about it before reading further.

\subsection{Transition Points}
The First Derivative Test is the only tool we need to distinguish a
minimum from a maximum. In the simplest case all we have to do is
identify each point where the tangent line is horizontal
($\dfdx{O}{x}=0$) and then determine if the slope of the tangent (the
derivative) is positive
($\dfdx{O}{x}>0$) or negative ($\dfdx{O}{x}<0$) on either side of that
point.



Clearly a local maximum occurs on the graph of $O(x)$ when the graph stops
increasing and starts decreasing. From Theorem~\ref{thm:FDT} we see
that if the derivative of $O(x)$ is positive then the function is
increasing, and when its derivative is negative then the function is
decreasing. Put this all together and it should be clear that to find
a maximum we need only find a place where the derivative \emph{transitions} from
positive to negative.

Similarly, to find a minimum we need to find a place where
the derivative \emph{transitions} from negative to positive. The most
obvious way for $\dfdx{y}{x}$ to make such a transition is for it to
pass through zero. We'll consider these first.

\begin{ProficiencyDrill}
  Can you think of another way to make that transition?
\end{ProficiencyDrill}

Let's check that proposition.
\begin{myexample}
  \label{problem:PPT}
  Let $y= O(x)=x^3-3x$ as above.

   The derivative of $O(x)$ is:
  $\dfdx{O}{x}=3x^2-3$, so the solutions of $\dfdx{O}{x}=0$ are
  $x=-1$, and $x=1$. We need to determine whether the derivative
  transitions from positive to negative, or from negative to positive at
  those two points.
  \begin{description}
  \item[At $x=-1$:] Pick a number less than $-1$, say $x=-2.$ Since
    $\dfdxat{O}{x}{-2}>0$ we see that $O(x)$ is increasing to the left
    of $x=-1.$ 
  \item[At $x=1$:] Similarly since $\dfdxat{O}{x}{2}<0$ we conclude that
    $O(x)$ is decreasing to the right of $x=1$.
  \end{description}
  Hold on though. We only tested one number on each side of our
  suspected extremum. Is that enough? Give this some thought. We will
  come back to it.
\end{myexample}


\begin{myexample}
  \label{example:x^5}
  Suppose $O(x)=x^5$. Find all extrema.

  Solving $\dfdxat{O}{x}{0}=5x^4=0$ we see that $x=0$ is the only
  possible extremum. But $5x^4>0$ for all $x$ so $\dfdx{O}{x}$ is
  positive on the left
  \emph{and on the right} of zero. So $O(x)=x^5$ does not have an
  extremum at $x=0$, despite the fact that $\dfdxat{O}{x}{0}=0.$
  this seems very curious until you graph the function.
  
  \begin{ProficiencyDrill}
    Sketch the graph of $O(x)$.       
  \end{ProficiencyDrill}
\end{myexample}

It should be clear that the solutions of $\dfdx{O}{x}=0$ are crucial
to solving optimization problems. This is because these are the only
places where
$O(x)$ can
transition from increasing-to-decreasing or from
decreasing-to-increasing. For that reason we will call these \term{possible
  transition points}. We solve optimization problems by locating
all possible transition points and identifying the extrema among
those. 

All solutions of $\dfdx{O}{x}=0$ are possible transition points but,
as we suggested earlier, there are others. We will hold off on a
formal definition until we have identified all types.

Returning to the question we asked at the end of
problem~\ref{problem:PPT} above, what do you think? Is it
enough to check only one point to determine the sign of the
derivative?

Of course it is. But \alert{only if we have identified all of the
  possible transition points.} Since the possible transition points
are precisely those places where $O(x)$ might change signs $O(x)$
\emph{must} either
be always positive or always negative between its between. We state
this formally as a lemma.

\begin{mylemma}
  \label{lemma:AdjRoots}
  Suppose $r_1$ and $r_2$ are two adjacent possible transition points of $f(x)$ (meaning
  that there is not another one between them). If $f(x)>0$ for some  $x$
  between $r_1$ and $r_2$, then $f(x)>0$ for every   $x$   between
  $r_1$ and $r_2$. Similarly, if $f(x)<0$ for some  $x$
  between $r_1$ and $r_2$, then $f(x)<0$ for every   $x$   between
  $r_1$ and $r_2$.
\end{mylemma}

\begin{ProficiencyDrill}
  Theorems and Lemmas in mathematics are always stated very formally
  and with as much generality as possible. There are good reasons for
  this but it does make it hard to read them sometimes.
  
  Read Lemma~\ref{lemma:AdjRoots} carefully and convince yourself that
  it really does say the same thing as the paragraph before it. 
\end{ProficiencyDrill}

\begin{embeddedproblem}
  \begin{wrapfigure}[]{r}{2in}
    \captionsetup{labelformat=empty}
    \centerline{\includegraphics*[height=1.4in,width=2in]{../Figures/ParabDist1}}
    \label{fig:ParabDist1}
  \end{wrapfigure}  In PIC \#\ref{problem:Parab1} we made the offhand remark, ``It is evident
  from the diagram below that there is no maximal square.''

  This is a misleading statement. 

  That is, it \emph{is} quite evident from the figure that there is no
  \emph{global} maximal square. But the sketch does not illuminate all
  of the subtleties of the problem. In fact this problem has three
  \emph{local} extrema: A global minimum (this is the one you were led to in
  PIC \#\ref{problem:Parab1}), a second local minimum, \emph{and} a local maximum.

  Find the other two extrema for this problem.

  The point here is this: diagrams and intuition are useful
  tools and you should use them whenever you can. But they have their
  limits and need to be supplemented. Diagrams can be badly drawn, and
  intuition can be misled.  They can't always be
  trusted. This is why you need to become proficient with the
  computational tools of Calculus.
\end{embeddedproblem}

\subsection{Undefined Derivatives}
\label{subsec:underl-deriv}
We said we would hold off on a formal definition of possible
transition points until we had identified the other kinds. But it hard
to see how there can be other kinds transition points isn't it? After
all, to be a
transition point is to have the derivative of the objective function
change signs; to \emph{transition}. Obviously, the only number that has negative numbers to
its left and positive numbers to its right (or \foreign{vice
  versa}) is zero, so don't the solutions of $\dfdx{O}{x}=0$ give us
\emph{all} of the possible transition points?

No, they  don't.

It is easy to jump to the conclusion that transition points \emph{only} occur
when $\dfdx{O}{x}=0$ but this is \emph{not} true. It is not even true
that  $\dfdx{O}{x}=0$ \emph{guarantees} a transition point  as we
saw in example~\ref{example:x^5}.



  It is the \emph{changing of the sign} of the derivative at a point
  that makes it a transition point, not having $\dfdx{O}{x}=0$.



\begin{myexample}
  Since it is the changing of the sign of the derivative that matters it
  is even possible for $x=a$ to be a transition point when
  $\dfdxat{O}{x}{a}$ is meaningless. (That is, when
  $\eval{\dfdx{O}{x}}{x}{0}$ is \underline{undefined}.

  \begin{ProficiencyDrill-1line}
  For example suppose
  $$
  O(x) = \sqrt{(1-x^{2/3})^{3}}.
  $$
  \begin{enumerate}[label={\bf{}(\alph*)}]
  \item Show that $\dfdx{O}{x}=-\frac{\sqrt{1-x^{2/3}}}{x^{1/3}}$, so
    that   $\eval{\dfdx{O}{x}}{x}{0}$ is undefined.
  \item     Verify that $\dfdx{O}{x}>0$ when $x<0$ and
    $\dfdx{O}{x}<0$ when $x>0$.
  \end{enumerate}
  \end{ProficiencyDrill-1line}
    From this example we can see that $\dfdx{O}{x}>0$ when $x<0$ and
    $\dfdx{O}{x}<0$ when $x>0$. Therefore $x=0$ is a transition point.
\end{myexample}


So we see that those points where $\dfdx{O}{x}=0$ and $\dfdx{O}{x}$ is
undefined are possible transition points.  There is one more kind.

\subsection{Optimization on a Closed Interval}
\label{sec:optim-clos-interv}
\markboth{\sc Optimization}{\sc Closed Interval} Real world
optimization problems will have a natural domain on which the
objective function is defined and it makes no sense to consider values
of $x$ outside of that domain. We actually saw this in
both Example~\ref{example:CircleDist} and
Example~\ref{example:HyperbolaDist1}. Recall that in  Example~\ref{example:CircleDist} we
needed to optimize the area of a square with one of its corners on the
unit circle: $x^2+y^2=1.$ Obviously it makes no sense to consider any
value of $x$ greater than $1$ or less than $-1$, so the natural domain
of this problem is $-1\le x\le 1$.
In Example~\ref{example:HyperbolaDist1} the natural domain was
everything outside the interval $(-1,1)$.
When that happens there is one more
kind of transition point to consider as we've seen.

\begin{myexample}
  \label{example:SqrClosed}
  \begin{wrapfigure}[]{r}{2in}
\captionsetup{labelformat=empty}
  \centerline{\includegraphics*[height=1in,width=2in]{../Figures/ObjSqr}}
\label{fig:}
\end{wrapfigure}
  Suppose that $O(x) = x^2$ and that the domain of the problem
  is $-1\le x\le 1$. 


  From the graph at the right we can see that this function has a (global)
  minimum at $(0,0)$ and  we can find this minimum by the methods we are
  already familiar with.

  We can also see that we have two (global) maxima. One at $(-1,1)$
  and the other at $(1,1)$.

  It should be clear what is going on here. If the domain were not
  restricted then the graph would continue to rise on both sides and
  there would be no maximum. But the restriction on the domain forces
  the graph to stop rising at the endpoints of the domain interval,
  making both of these endpoints local maxima. When the domain is thus
  restricted there is the possibility that a maximum (or minimum)
  might occur at the endpoints of the domain interval as well. As a
  result we must always check the endpoints of the domain interval (if
  there are any) when searching for extrema.

  We have located all of
  the places where extrema can occur so we will define them
  formally.

\end{myexample}
\begin{mydefinition}[Possible Transition Points]
  \label{def:TransitionPoints}
  The \term{possible transition points} of an objective function, $O(x)$, come in three
  categories:
  \begin{enumerate}
  \item Any point, $x=a$, where $\dfdxat{O}{x}{a}=0$ is a possible transition
    point.
  \item Any point, $x=a$, where $\dfdxat{O}{x}{a}$ is not defined is
    a possible transition point.
  \item If the domain of $O(x)$ is \alert{closed}, then the endpoints
    of any intervals in the domain are possible transition
    points\aside{Since no ``transition'' occurs at an endpoint we
      could object that it is incorrect to call the endpoints
      ``transition points.'' But since the domain is the whole univers
      for the problem we could say that the graph transitions from
      non-existence to existence at an endpoint. In any case our
      interest is in finding all possible extrema, not in adhering to
      the strictures of the English language.}.
  \end{enumerate}
\end{mydefinition}

Does it seem odd to you that only two of the three conditions above
involve the derivative? Strictly speaking, there are really only two
conditions because, as we will see, the derivative is also undefined
at the endpoints of closed intervals. However this is a technicality
which doesn't affect anything we are doing now so we won't fuss about
it. We will return to this in Section~\ref{sec:one-sided-deriv}.

As we said, real world problems usually come with a natural
domain. But we're not looking at real world problems in this section.
The upshot of this is that to explore the issues we're interested in
we will have to artificially impose constraints on the the domain of a
given problem.

For example, in Example~\ref{example:SqrClosed} we artificially
specified a domain of $[-1,1]$. When a domain is imposed artificially
like this it is easy to think of it as ``not really there.'' This is a
mistake. The specification of the domain is a part of the problem,
just like the specification of the function. Think of the domain
specification at the ``universe'' for that particular problem. For
Example~\ref{example:SqrClosed} nothing outside of the interval
$[-1,1]$ exists.


So to optimize a function all we need to do is find the possible
transition points, evaluate our objective function at these points,
and determine which is largest and which is smallest, right?

Well, yes and no.

If the domain of our problem consists of one or more closed intervals, then
yes. That's all there is to it. Really. It is that
straightforward\aside{Remember that ``straightforward'' does not
  mean ``easy.''}. The following theorem summarizes the
optimization procedure when the domain is a closed interval.

\begin{mytheorem} [Max/Min Theorem]
  \label{thm:max-min}
  If $f(x)$ is a continuous function defined on a closed interval
  $[a,b]$ then there is at least one point in $[a,b]$, say
  $c$, such that $f(c)$ is the
  maximum value of the function on the interval. Also there is also at least
  one point, say $d,$ in $[a,b]$ such that $f(d)$ is the
  minimum. Moreover $c$ and $d$ will either be endpoints of the interval,
  places where   $\left.\dfdx{f}{x}\right|_{x=c} = 0 =
  \left.\dfdx{f}{x}\right|_{x=d},$ or where the derivative does not exist.
\end{mytheorem}
If the domain is all real numbers then things get a little more
complicated. We will address this situation in the next section.

\begin{ProficiencyDrill}
  Read the Max/Min Theorem carefully and convince yourself that it
  says what we said it says.
\end{ProficiencyDrill}

  \begin{embeddedproblem}{}
    \begin{enumerate}[label={\bf{}(\alph*)}]
    \item Notice that $y=x^2$ has no extrema on $(0,\infty).$ Explain
      why this does not violate the Max/Min Theorem?
    \item Notice that $y=\frac{1}{x}$ has no extrema on $[-1,1].$
      Explain why this does not violate the Max/Min Theorem?
    \end{enumerate}
  \end{embeddedproblem}

We've actually been using Theorem~\ref{thm:max-min} in an informal way
all along.  In Example~\ref{example:CircleDist} and Proficiency
Drill~\ref{problem:CircleDist2} we followed essentially the procedure
indicated by Theorem~\ref{thm:max-min}.



\begin{ProblemSection}

  \begin{embeddedproblem}{}
    Find the absolute maximum and absolute minimum of the function
    $$
    f(x) = x^3-3x,
    $$
    on each of the following intervals.
    Verify your solution with   graphing software.
    \begin{enumerate}[label={\bf{}(\alph*)}]
      \begin{multicols}{2}
      \item $[2 , 5]$ \solution{Minimum:  $(2,2)$,  Maximum:
          $(5,110)$}
      \item $[1,2]$ \solution{Minimum:  $(1,-2)$,  Maximum:
          $(2,2)$}
      \item $[0,1]$ \solution{Minimum:  $(1,-2)$,  Maximum:
          $(0,0)$}
      \item $[-1,0]$ \solution{Minimum:  $(0,0)$,  Maximum:
          $(-1,2)$}
      \item $[-2,-1]$ \solution{Minimum:  $(-2,-2)$,  Maximum:
          $(-1,2)$}
      \item $[-\infty,-2] $ \solution{Minimum:  $(-1,2)$,  Maximum:
          $\displaystyle  y=$}
      \end{multicols}
    \end{enumerate}

  \end{embeddedproblem}

  \begin{embeddedproblem}{}
    Find the absolute maximum and absolute minimum of each of the
    following functions on the given intervals.  on each of the
    following intervals.  Verify your solution with graphing software.
    \begin{enumerate}[label={\bf{}(\alph*)}]
      \begin{multicols}{2}
      \item $  y= x^2$ \\ on the interval $[-5 ,5 ]$ \solution{          Minimum:
          $(0,0) $           Maximum:  $(-5,25) $ and
          $(5,25) $}
      \item $  y= x^3$ \\ on the interval $[-5 ,5 ]$ \solution{          Minimum:
          $(-5, -125) $           Maximum:  $(5, 125) $}
      \item $   y= (x-1)(x+2)$ \\ on the interval $[0 ,2 ]$ \solution{
          Minimum:  $(0,-2) $               Maximum:
          $(2,4) $}
      \item $   y= x(x-1)(x+2)$ \\ on the interval $[2, 5]$ \solution{
          Minimum:  $(2, 8) $               Maximum:  $(5,
          140) $}
      \item $   y= x(x-1)(x+2)$ \\ on the interval $[0 ,2 ]$ \solution{
          Minimum:  $\left(\frac{-1+\sqrt{7}}{3}\right)$
          Maximum:  $(2,8) $}
      \item $ y = e^{-x^2}$ \\ on the interval $[2 ,10 ]$ \solution{
          Minimum:  $(10,e^{-100}) $
          Maximum:  $(2, e^{-4}) $}
      \item $ y = e^{-x^2}$ \\ on the interval $[-1 ,2 ]$ \solution{
          Minimum:  $(-1,e) $                 Maximum:
          $(2, e^{-4}) $}
      \item $ y = e^{-\frac{1}{x^2}}$ \\ on the interval $[2 ,10 ]$ \solution{
          Minimum:  $(2, e^{-1/4}) $
          Maximum:  $(10, w^{-1/100}) $}
      \item $ y =(x^3-27)^{\frac13}$ \\ on the interval $[-5 , 5]$ \solution{
          Minimum:  $(-5,\sqrt[3]{-152}) $
          Maximum:  $(-5,\sqrt[3]{98}) $}
      \end{multicols}
    \end{enumerate}
  \end{embeddedproblem}

  \begin{embeddedproblem}{}
    Find the absolute maximum\notetoself{Add domains that consist of
      multiple closed intervals.  -- Jan 10, 2021} and absolute minimum of each of the
    following functions on the given intervals.  on each of the
    following intervals.  Verify your solution with graphing software.
    % <figure>
    % 
    % <interactive desmos="xlit0cc7vm"
    % width="150%" aspect="2:1" />
    % 
    % </figure>
    \begin{enumerate}[label={\bf{}(\alph*)}]
      \begin{multicols}{2}
      \item $\displaystyle y=\frac{x(x^2-1)e^x}{x^3-1}$,\\ on the interval $[-2, 0]$
      \item $\displaystyle y =x^2\sin(x^2)$, \\ on the interval$[0 ,2 ]$
      \item $\displaystyle y =\frac{\ln(x)}{x}$, \\ on the interval $[2 ,10 ]$
      \item $\displaystyle y = \frac1x$, \\ on the interval $[1/2 ,4 ]$
      \item $\displaystyle y =\abs{x}$, \\ on the interval $[-1 ,1 ]$
      \item $\displaystyle y =\abs{x-2}$, \\ on the interval $[-4,4 ]$
      \item $\displaystyle y=\frac{7}{x^2+5}$, \\ on the interval $[-2,2]$
      \item $\displaystyle y=4-x^3$, \\ on the interval $[-2,1]$
      \item $\displaystyle y=-\frac1x$, \\ on the interval $[1/2,3]$
      \item $\displaystyle y=\sqrt[3]{x}$, \\ on the interval $[-1,8]$
      \item $\displaystyle y=-3x^{2/3}$, \\ on the interval $[-1,1]$
      \item $\displaystyle y=\sec\theta$, \\ on the interval $[-\pi/3,\pi/6]$
      \item $\displaystyle y=\frac1x+\ln(x)$, \\ on the interval $[1/2, 4]$
      \item $\displaystyle y=e^{-x^2}$, \\ on the interval $[-2, 1]$
      \item $\displaystyle y=\frac{x}{x^2-x+1}$, \\ on the interval $[0,3]$
      \item $\displaystyle y=2\cos x+\sin 2x$, \\ on the interval $[0,\pi/2]$
      \item $\displaystyle y= 2x+\frac{1}{2x} $ \\ on the interval $ [1,4]$
      \item $\displaystyle y=\frac{1-x}{x^2+3}$ \\ on the interval $[-2,5]$
      \item $\displaystyle y=x\sqrt{1-x^2} $ \\  on the interval $[-1,1]$
      \item $\displaystyle y=x\sqrt{4-x^2} $ \\ on the interval $[0,2]$
      \item $\displaystyle y=x(2-x)^{1/3}$ \\ on the interval $[1,3]$
      \item $\displaystyle x^{1/2}-x^{3/2}$ \\ on the interval $[0,4]$
      \end{multicols}
    \end{enumerate}
  \end{embeddedproblem}
  
  \begin{embeddedproblem}{}
    Find the maximum and minimum of the function
    $$
    f(x) = x^3-12x
    $$
    on each of the following intervals.
    \begin{enumerate}[label={\bf{}(\alph*)}]{}
      \begin{multicols}{3}
      \item $[0,4]$
      \item $[-4,0]$
      \item $[-2,2]$
      \item $[-1,3]$        
      \item $[-2, 1]$
      \item $[-\infty,\infty]$
      \end{multicols}
    \end{enumerate}
  \end{embeddedproblem}

  \begin{embeddedproblem}{}
    For each of the following functions defined on specific closed,
    bounded intervals, find the maximum and minimum of the function on
    that interval.
    \begin{enumerate}[label={\bf{}(\alph*)}]
      \begin{multicols}{2}
      \item $f(x)=2x^3+3x^2-12x+1  \text{ on } [-3,2]$
      \item $f(t)=2t^3+3t^2-12t+1   \text{ on } [-3,3]$
      \item $g(x)=4x^{\frac{3}{2}}-3x^2-2  \text{ on }[0,2]$
      \item $h(z)=3z^(2/3)+2z-4  \text{ on }[-2,1]$
      \item $L(x)=\ln((x-2)^6) -x^2  \text{ on }[-1.5,4]$
      \item $f(\theta)=1/2 \theta-\sin\theta \text{  on }[0,2\pi]$
      \end{multicols}
    \end{enumerate}
  \end{embeddedproblem}

  \begin{ProficiencyDrill}
    Consider the function  $ f(x) = x^3-12x$.
    \begin{enumerate}[label={\bf{}(\alph*)}]
    \item Show that on the interval $[-4, 4]$ $f(x)$
      attains a maximum at both $x=-2$ and $x=4$ and
      a minimum at both $x=2$ and $x=-4$.
    \item Find the extrema of $f(x)$ on the interval
      $[-3,3]$.
    \end{enumerate}
  \end{ProficiencyDrill}
    
  \begin{embeddedproblem}{}
    Two positive numbers sum to $10.$  Find the maximum product of these two numbers.
  \end{embeddedproblem}

  \begin{embeddedproblem}{}
    The product of two positive numbers is $10.$  Find the minimum sum
    of these two numbers.  Would there be a maximum sum?  Explain.
  \end{embeddedproblem}


  \begin{embeddedproblem}{}
    Suppose we want to build a rectangular corral using $2000$ feet of
    fencing.
    \begin{enumerate}[label={\bf{}(\alph*)}]
    \item Find the dimensions of the corral with the maximum possible
      area, if there is one.
    \item Find the dimensions of the corral with the minimum possible
      area, if there is one.
    \end{enumerate}
    Does this problem seem familiar to you?
  \end{embeddedproblem}

  \begin{embeddedproblem}[(Variations on a theme)]
    You\notetoself{I made these problems up on the fly. I have not
      yet checked to see if there are good problems, or even if they
      are solveable.} have a piece of wire $100$ inches long.
    \begin{enumerate}[label={\bf{}(\alph*)}]
    \item You cut the wire into two pieces.  Find the length of each
      piece of the wire that yields the largest area bounded by the
      two pieces, given the following constraints.
      \begin{enumerate}[label={\bf{}(\roman*)}]
      \item Each piece is to be bent into a square.
      \item Each piece is to be bent into a circle.
      \item The first piece is to be bent into a square and the second
        into a circle. 
      \item The first piece is to be bent into a rectangle with one
        side three times as long as the other.  the second piece is to
        be bent into a square.
      \end{enumerate}
    \item You cut the wire into three pieces.  Find the length of each
      piece of the wire\notetoself{Need to add more variations to this
        problem.} that yields the largest area bounded by the three
      pieces, given the following constraints.
      \begin{enumerate}[label={\bf{}(\roman*)}]
      \item Each piece is to be bent into a square. The length of the sides of
        second square are twice the length of the sides of the
        first. The length of the sides of the third are three times
        the length of the second.
      \item Each piece is to be bent into a circle. The radius of the
        second circle is twice the radius of the first. The radius of
        the third is three times the radius of the second.
      \item Two pieces are to be bent into a square and the third into
        a circle.
      \end{enumerate}
    \end{enumerate}
  \end{embeddedproblem}

  % \begin{embeddedproblem}{}
  %   Let $(x,y)$ be a generic point on the unit circle $x^2+y^2=1.$  Find the
  %   maximum and minimum values of the slope of the line joining $(x,y)$ with
  %   $(4,-3).$\\
  %   \centerline{\includegraphics*[height=1.6in,width=2.2in]{../Figures/Lagrange1}}
  %   From the sketch above it should be clear that Of all
  %   lines through the point $(4, -3)$ that also intersect the circle the
  %   ones  we seek are have the maximum and minimum possible slopes.
  % \end{embeddedproblem}

  \begin{embeddedproblem}{}
    Consider the side view of a person at point $P$ looking at a
    billboard sign in the diagram below\\
    \centerline{\includegraphics*[height=2in,width=2in]{../Figures/BillBoardProblem}}
    \begin{enumerate}[label={\bf{}(\alph*)}]
    \item Find a formula for the viewing angle $A$ in terms of $x.$ 
      Don't forget to put in the range for possible values for $x.$
    \item Find the value of $x$ that will maximize the viewing angle $A.$ 
    \end{enumerate}
  \end{embeddedproblem}


  \begin{embeddedproblem}
    Find the value of $r$ that maximizes $A(r)$ in the following
    diagram:\\
    \centerline{\includegraphics*[height=1.5in,width=1.5in]{../Figures/ArchimedPump1}}
  \end{embeddedproblem}

  \begin{embeddedproblem-enumerate}[(Variations on a theme)]
  %   \begin{enumerate}[label={\bf{}(\Roman*)}]
  %     \item\  \\
  %   \begin{enumerate}[label={\bf{}(\alph*)}]
  %   \item The strength of a rectangular beam is proportional to the
  %     product of width $w$ of the beam and the square of the height $h$
  %     of the beam.  Find the dimensions of the strongest rectangular
  %     beam which can
  %     cut from the following circular log.\\
  %     \centerline{\includegraphics*[height=2in,width=2in]{../Figures/BeamProblem}}
  %   \item The stiffness of a rectangular beam is proportional to the
  %     product of width $w$ of the beam and the cube of the height $h$ of
  %     the beam. Find the dimensions of the stiffest possible beam.
  %   \end{enumerate}
  % \item 
    \begin{enumerate}[label={\bf{}(\alph*)}]
    \item The \term{strength} of a rectangular wooden
      beam is proportional to its width times the square of its
      height.  What is the ratio of the height to width of the
      strongest beam that can be cut from a cylindrical log?
    \item The \term{stiffness} of a rectangular wooden beam is proportional
      to its width times the cube of its height.  What is the ratio of
      the height to width of the stiffest beam that can be cut from a
      cylindrical log?  Compare this to the answer in the previous
      problem.
    \item Suppose we have an  objective function -- let's call it
      the \emph{generalized stiffness} -- which is proportional
      to the width times the $n$th power of the height? What would the
      optimal ratio be now?
    \end{enumerate}
%   \end{enumerate}
  \end{embeddedproblem-enumerate}

  As we’ve said before, Calculus is merely one tool in your problem
  solving toolbox.  You will need all of your mathematical experience
  to transform a problem into a form which allows you to bring your
  Calculus to bear on it. Moreover there is typically a lot of algebra
  to perform after Calculus is used.  Quite often it is a ``zero sum
  game''; simplifying as much as possible with algebra, etc. leads to
  easier calculus computations. Conversely less non-Calculus work
  leads to more challenging differentiation.  With experience comes
  the ability to balance the workload across your skills. When you
  realize that one approach will lead to a complication, it can be 
  beneficial to look for another approach.  The following example shows
  that different approaches to the same problem can lead to different
  challenges.
  \begin{myexample}
    \label{example:RodAroundACorner}
    Consider the following top view of a pole of length $P$ being
    rotated horizontally around the $90$ degree corner of two hallways
    of widths $a$, and $b$.  The problem is to find the longest pole
    which can be turned around the corner.\\
    \centerline{\includegraphics*[height=2in,width=5in]{../Figures/RodRoundCorner1}}
    \begin{description}
    \item [First Approach:] Let’s take the most direct approach by
      labeling a number of quantities and writing down our objective
      function and constraints.  With this in mind, consider the
      following labeling.\\
      \centerline{\includegraphics*[height=2in,width=5in]{../Figures/RodRoundCorner2}}
      \begin{embeddedproblem-enumerate}
        \begin{enumerate}[label={\bf{}(\alph*)}]
        \item Show that with the above labelling we have the
          objective: function $L^2=x^2+y^2$ and the constraint:
          $\frac{x-a}{b}=\frac{a}{y-b}$.
        \item Show that differentiating the objective function and
          constraint leads to having to solve the following quartic
          equation
          $$
          b^2(b+ab(x-a))=a^2x(x-a)^3.
          $$
        \end{enumerate}
      \end{embeddedproblem-enumerate}
        There is a quartic formula, but it is pretty unwieldy.  So, let’s try another approach.  
    \item [Second Approach:] Consider the same problem formulated in a
    different way.  We can put in a set of axes and plot the point
    $(a,b)$.  The problem now is to look at all the family of all
    lines with negative slope, $m$, passing through point $(a,b)$ and
    find the value of $m$ which minimizes the distance $L$ between the
    $x$ and $y$ intercepts.\\
    \centerline{\includegraphics*[height=1in,width=2in]{../Figures/RodRoundCorner3}}
    \begin{embeddedproblem-enumerate}
      \begin{enumerate}[label={\bf{}(\alph*)}]
        \item With this approach, what would our objective function
          $L$ be?
        \item As to constraints, find $z$ and $w$ in terms of the
          slope $m$  and show
          that using this approach leads to solving the quartic
          equation
          $$
          0=-b^2+abm-abm^3+a^2m^4.
          $$
        \end{enumerate}
    \end{embeddedproblem-enumerate}
        Again, this is doable, but not very appealing.  Let’s look at
        another approach.
  \item [Third Approach:]  Consider the following labeling.\\
    \centerline{\includegraphics*[height=1in,width=2in]{../Figures/RodRoundCorner4}}
    \begin{enumerate}[label={\bf{}(\alph*)}]
      \item Show that $L$ is minimized when
        $\tan(\alpha)=\sqrt[3]{\frac{b}{a}}$.
      \item Use the value obtained in part a to show that the minimum
        value of $L$ is $\left(a^\frac23+b^\frac23\right)^\frac32$.
      \item What would the minimum of $L$  be in the special case
        where $a=b$?  Does this answer make sense geometrically? 
      \end{enumerate}
    \end{description}
    Notice that in the third approach, the constraints required a bit
    more difficult differentiation, but led to a equation that we
    could more readily solve.
  \end{myexample}
  
  \begin{embeddedproblem}
    The illumination of an object by a light source is proportional to
    the strength of the light source and inversely proportion to the
    square of the distance between the source and the object.
    Consider the following diagram of two light bulbs of equal
    strength a distance of $D$ apart and an object at point $P$ on a line
    segment parallel to and distance $E$ from the line segment joining
    the two bulbs. \\
    \centerline{\includegraphics*[height=1.5in,width=4.5in]{../Figures/LightBulbProblem}}
    \begin{enumerate}[label={\bf{}(\alph*)}]
    \item Find a formula for the total illumination of the two light
      bulbs on the point $P$ in terms of $x,$ $D,$ and $E.$  (Don't forget that
      $0\le x\le D.$)
    \item Show that $I$ has a local minimum for $x=D/2.$  (This
      should not come as a surprise.)
    \item Show that for $E>\frac{D}{\sqrt{2}},$ $x=D/2$
      (surprisingly?) does not yield the absolute minimum. (This shows
      the importance of checking endpoints.)
    \end{enumerate}
  \end{embeddedproblem}



  \begin{embeddedproblem}{}
    For each of the IVPs below, answer these questions:
    \begin{enumerate}
    \item On what intervals is $y(t)$ increasing and on what
      intervals is it decreasing?  Identify all local extrema.
    \item On what intervals is the function concave up and on what
      intervals is it concave down?  Identify all inflection points.
    \item Plot a reasonable graph for $y.$ 
    \end{enumerate}
    \begin{enumerate}[label={\bf{}(\alph*)}]
    \item $\dfdx{y}{t}=\sqrt{t},\ \         y(0)=1$
    \item $\dfdx{y}{t}=x^3-x,\ \      y(1)=3$
    \item $\dfdx{y}{t}=x-x^3,\ \       y(1)=3$   \hint{Compare this to
        the previous problem?}
    \end{enumerate}
  \end{embeddedproblem}
\end{ProblemSection}

\subsection{Optimization on a Non-Closed Interval}
\label{sec:optim-non-clos}

Theorem~\ref{thm:max-min} explicitly assumed that we are on a
\underline{closed} interval. This assumption was crucial. Without it
even simple functions that you would expect to have extrema do not. 
\begin{myexample}
\label{example:OpenNoExtrema}
  For example, suppose our objective function is $O(x)=x$, and that
  the domain is the interval $(0,1)$ so the interval is not
  closed. (Recall that the notation $(0,1)$ means that the endpoints
  are \emph{not} included.)

  Before we go on give this some thought. Do you see why the graph of
  this function has neither a minimum nor a maximum?
  If not here is one way to prove it.

  \begin{myproof}
    If $x$ is any number between $-1$ and $1$ then $f(x)$ can't be a
    maximum (minimum) because any number just to the right (left) of $x$
    will give us a $y$-coordinate that is higher (lower) than $f(x)$.

    How do we know this? Because the derivative of $O(x)$ is $1$, which
    says that $O(x)$ is always increasing.

    Can we always find a number ``just to the right'' of $x$ without
    leaving the interval $(-1,1)?$ Sure. Just take the number halfway
    between $x$ and $1.$
  \end{myproof}
\end{myexample}

This leaves us with a dilemma: If the domain is not closed we can't
rely on the Min/Max Theorem to guarantee the existence of extrema. In
fact, in this situation Example~\ref{example:OpenNoExtrema} shows that
there might not even be any extrema. But if there are we'd still like
to find them. 

This is a problem but it is not insurmountable. In fact we already
have all of the tools we need. We just need to use them more
carefully. It is still true that local extrema can only occur at
transition points, so we still have to find all of the possible
transition points.

An example will make this clearer.
\begin{myexample}
  Before we start we admonish you to \alert{not} graph this
  example until after you have understood the analysis of the
  problem. The goal of this section is for you to learn to analyze a
  problem using your Calculus tools. This works best if you use
  Calculus to visualize an unknown graph, rather than to confirm a
  known graph.

    Suppose $O(x)=3x^3-9x$ is our objective
    function and the domain is all real numbers ($\RR$). Locate all
    possible points of transition and determine any global extrema, if
    such exist.

    We leave it to you to show that the possible points of transition
    are $x=-1$ and $x=1$. Thus the intervals between them are
    $(-\infty,-1)$, $(-1,1)$, and
    $( 1,\infty)$.
    \begin{center}
      \begin{minipage}{.75\linewidth}
        \begin{mynotation}{``Infinity'' ($\infty$)} For convenience,
          and to be consistent with the interval notation, we adopt
          the convention that $(1,\infty)$ denotes the positive real
          numbers greater than $1$.

          Don't let this make you crazy.  Clearly there is not a
          number at the right endpoint of the number line called
          ``infinity.''  There can't be because the number line
          doesn't have a right endpoint. Nor does it have a left
          endpoint. The notation $(1,\infty),$ is just a mnemonic
          convention that means ``the interval which includes all
          numbers to the right of one.'' In similar fashion the
          notation $(-\infty,-1)$ means the interval to the left of
          one. Obviously, then $(-\infty,\infty)$ means the entire
          real number line: $(-\infty, \infty)= \RR$.
        \end{mynotation}
      \end{minipage}
    \end{center}
    It is easy to show that\\
    \centerline{\hfill{}$\dfdxlt{O}{x}{-1}>0$,\hfill{}
      $\dfdxlgt{O}{x}{-1}{1}<0$, \hfill{}and
      \hfill{}$\dfdxgt{O}{x}{1}>0$.\hfill}

    We conclude that $O(x)$ is increasing on $(-\infty,-1)$,
    decreasing on $(-1,1)$, and increasing again on $(1,\infty)$.  Can
    we conclude that we have a global maximum at $x=-1$ and a global
    minimum at $x=1$?

    No we can't because this problem has no global extrema.

    Now plot a graph of the function. Clearly we have no global
    extrema at all because on the left as $x$ gets larger \emph{in the
      negative direction} so does $O(x)$ and on the right as $x$ gets
    larger \emph{in the positive direction} so does $O(x)$.


    This is awkward to say, so we usually say something like\aside{Do
      not try to find too much meaning in these phrases. They clearly
      can not possibly mean what they pretend to mean. $x$ is a number
      but infinity is not, so it is meaningless to say that $x$ ``goes
      to'' infinity.} ``as $x$ goes infinite, $O(x)$ also goes
    infinite'' and ``as $x$ goes negatively infinite, $O(x)$ also goes
    negatively infinite.''
  \end{myexample}
  \begin{myexample}
    Now suppose that the objective function is
    $O(x)=\frac{-x}{x^2+1}$, and suppose we want to find any global
    extrema. We proceed as before.

    Differentiating we have
    $\dfdx{O}{x}=\frac{x^2-1}{(1+x^2)^2}$. Setting $\dfdx{O}{x}=0$ and
    solving we see that we have the same possible transition points as
    before: $x=-1$, and $x=1$.

    Also, just like the previous example our objective function
    increases from $-\infty$ to $-1$, decreases from $-1$ to $1$ and
    then increases again from $1$ to $\infty$.  This time we actually
    do have a global maxmimum at $x=-1$ and a global minimum at
    $x=1$. Plot the graph to see. 
\end{myexample}

  
  In neither of the two previous examples can we conclude that we have
  global extrema \emph{based only on our analysis of the derivative}
  because the derivative gave us the same information in both
  examples. In order to tell the two situations apart we needed more
  information. For these examples we took that information from our
  graphs. This is ok for now but eventually we'll need some analytical
  method. We'll come back to this is Chapter~\ref{chapt:shape-things}.

  \begin{Digression}[What Does ``Increasing'' Mean?]

In\notetoself{I'm beginning to think this whole section needs to
  disappear.  -- Dec 30, 2020} the previous sections we saw that the derivative gives us
information about curves in much the same way that the slope of a
straight line gives us information about the line. The difference
is that since the numerical value of the derivative changes at
each point it only gives us \emph{local} information about the
curve at that point. Whereas slope gives us global information
about the whole line.

As simple and intuitive as this seems to be there are two
problems.

The first is, what do we mean by increasing and decreasing?

No, really. What do we mean? Try to give a single all-encompassing
definition of ``increasing'' that includes all of the following:
\begin{enumerate}
\item A variable, say $x$, is increasing as it moves from left to
  right on the horizontal axis.
\item A variable, say $y$, is increasing as it moves from lower to
  higher  on the vertical axis. 
\item In the northern hemisphere the average daily temperature tends
  to increase from January to July.
\item The pitch of a whistle becomes more shrill as its frequency
  increases.
\item As the color of a light becomes bluer as its frequency increases.
\item When you eat ice cream your pleasure increases.
\end{enumerate}

\begin{wrapfigure}[]{r}{1.5in}
  \captionsetup{labelformat=empty}
  \centerline{\includegraphics*[height=1.5in,width=1.5in]{../Figures/Abstract1}}
  % \label{fig:Abstract1}
\end{wrapfigure}
In the context of mathematics -- specifically when looking at the graph
of a function -- you've been taught all of your life that increasing
means ``rising to the right'' and that decreasing means ``dropping to
the right.'' Thus the graph at the right is usually said to be is
increasing from the far left to a point somewhere between zero and
one, decreasing from that point to a point somewhere between three and
four, and then increasing again thereafter.  That interpretation is
deeply tied to the convention of saying that as we go farther to the
right on the $x$-axis the values of $x$ get bigger. But this is only a
convention.

Conventions are handy, they are even necessary, but they are not
reality. If for some reason we had adopted the convention that
moving farther to the \emph{left} causess the value of $x$ to increase
then we would say (reading right-to-left) that this graph decreases
first, then increases for a while, and then decreases again.

The good news here is that  we will continue to adopt the usual
convention. The value of that convention is that when
we say a graph is increasing or decreasing we will be understood
without having to first identify which direction is which.

\alert{But} keep in mind that what this really means is that the
$y$ coordinate of the graph of our function, $f(x),$ becomes
numerically larger  as $x$ becomes numerically larger. Like we said, conventions are handy, but we can't let them
rule us. Which direction means ``bigger'' depends entirely on the
coordinate system we choose for a given problem, and we 
get to choose the coordinate system that we like. We are not slaves to
our own conventions.

Resolving the meaning of ``increasing'' was simply a matter of recognizing our own
convention.  Difficulties arising from too much reliance on a
convention are usually easy to resolve once we realize that's what
they are.  However, distinguishing between a convention (especially
one as old and comfortable as this one) and reality can be very
difficult because conventions, once we've gotten comfortable with
them, begin to feel fundamental; as if they \alert{are} reality. They are not.

But we have another problem and this one cannot be resolved by
noticing that we've mistaken a convention for reality. We stated above
that if the derivative of some function $f(x),$ is positive at some
point $x=a,$ then $f(x)$ is increasing \emph{at that point.}

The problem is this: What does it mean to be increasing at a single point?

In order to know that $f(x)$ has increased we need to move away from
$x=a$ to some point (to the right of $a$ in our standard convention),
say $x=a+h$, where $h$ is positive, and see if $f(a+h)$ is bigger than
$f(a).$ But even if $f(a+h)$ \emph{is} bigger the only information we
have is that $f(a+h)>f(a).$ This does not guarantee that $f(x)$ was
increasing at every point between $a$ and $a+h$. There might
be some point between $a$ and $a+h$ -- let's call it $\alpha,$ --
where $f(\alpha)<f(a).$ Even worse there might be some point, $\beta$
between $a$ and $\alpha$ where $f(\beta)>f(a).$ Yet again there might
\emph{another} point $\gamma$ between $a$ and $\beta$ where
$f(\gamma) < f(a),$ again. And so on $\ldots$

If the distance between $a$ and $a+h$ is very small (that is, if $h$
is very small) and if $f(a+h)>f(a),$ it is very easy to \emph{believe}
that $f(x)$ increased through all of the points from $a$ to $a+h.$ But
belief is not knowledge. In fact, we have no information whatsoever
about the behavior of $f(x)$ between $a$ and $a+h.$ We simply know
that when $x=a+h$ the function was greater than it was at $x=a.$

This is getting very hard to think about isn't it? That's because what
we are getting caught on is one of Zeno's Paradoxes. Specifically, the
arrow paradox. Briefly, Zeno\aside{If you have never heard of the
  Greek philosopher Zeno, you really should \google{} him. He's an
  interesting guy.  But wait until you are done with your Calculus for
  today.} said that an arrow in flight can't really move. This is
because to get from where it is now, to any other point in space it
has to traverse all of the (infinitely many) points in between. But
this is also true of all of those points. Thus before the arrow can
move once it has to move infinitely many times. That's just crazy, so
Zeno concluded that the arrow must not move at all.

Here's how we'll get around this difficulty. We all know what we mean
when we say that a curve is increasing at a point. The mental image
that is generated is very clear. The problem is that any attempt to
put this into words runs headlong into Zeno and his paradox. We cannot
resolve the paradox, so we will simply go around it for now. That is,
we first recognize that the usual intuitive meaning of the word
increasing does not apply at a single point, since we need at least
two points to say we've increased from one to the other. Thus the
phrase ``increasing at a point'' is without meaning.



If the phrase has no meaning then we can \emph{assign} meaning to it
in any manner we choose. That's what a definition is, an assignment of
meaning. We choose
the following definition. It will be very familiar:\\
\begin{mydefinition}[Increasing and Decreasinging at a Point]
  \label{def:first-deriv-test}
  Consider the graph of the function $y(x),$ at some point
  $(a, y(a)).$
  \begin{enumerate}
  \item The graph is said to be \emph{increasing} at $x=a$ if
    $\dfdxat{y}{x}{x=a} > 0$
  \item The graph is said to be \emph{decreasing}\footnote{Notice that
      our definition does \underline{\bf{}not} address the situation
      $\dfdxat{y}{x}{a}=0.$ Since the definition does not address
      this we say that the function is neither increasing nor
      decreasing when $\dfdxat{y}{x}{a}=0.$ } at $x=a$ if
    $\dfdxat{y}{x}{a} < 0$
  \end{enumerate}
\end{mydefinition}

There are two difficulties with this definition. The first is that it
probably feels like we've taken something simple and made it
complicated for no good reason. After all, what have we really
accomplished? We took the highly intuitive observation that we called
the First Derivative Test and we turned it around to make it a
definition. It probably seems like we're blathering on about nothing
at all, making easy things seem difficult and getting lost in a jungle
of logical pecadillos and notations. This is the usual criticism
brought against mathematicians and mathematics by people who haven't
studied much mathematics.

But stop and think about this just a little more deeply. Is that
really all that we've done?

Suppose we would like to know when a function is increasing and when
it is decreasing (we would). In the absence of any other tools, our
only option is to look at the value of the function at each point
individually and decide whether or not it is increasing at each
point. But if we are considering them one at a time the statement
``increasing at a point'' needs to be meaningful. That's what
Definition~\ref{def:first-deriv-test} does for us. Moreover it gives us exactly the meaning we need for the First
Derivative Test. 

If $\dfdxat{f}{x}{a}>0$ then \alert{by definition} $f$ is increasing
at $a$ and this definition matches our intuitive understanding of
``increasing'' precisely. It must.
The other difficulty is that this definition highlights an apparent
discrepancy between the dynamic and static views of Calculus.

When we take the dynamic (Newtonian) viewpoint there seems to be no
problem with using Definition~\ref{def:first-deriv-test}. If $p(t)$
represents the position of an object in motion,
$\dfdxat{p}{t}{\tau}>0$  means that at time $\tau$ the point is
moving forward. All of the intuition built over a lifetime of living
in this world tells us that if an object is moving forward at time
$\tau$ then it will keep moving forward, at least for a little while
(that is, on some interval), after time $\tau$.  We call this
``inertia''.

But when we take the static (Leibnizian) viewpoint there is a
problem. In this case $\dfdxat{p}{t}{\tau}>0$ means only that the
slope of the graph of $p(t)$ is positive. Absolutely nothing about the
geometry of the graph or its tangent line forces $p(t)$ to continue
increasing (have a positive derivative) to the right of $\tau$.

This discrepancy shines a light on the limitations of
Definition~\ref{def:first-deriv-test}. If we use it to show that
$p(t)$ is increasing \emph{at} $\tau$ all we really know is that:
\begin{enumerate}
\item Its velocity at $\tau$ is positive (dynamical viewpoint)
  \alert{at $\tau$}, or that,
\item Its slope is positive (static viewpoint) \alert{at $\tau$}.
\end{enumerate}
Neither of these actually tells
us anything about what happens  \alert{after} $\tau$.

To be sure, we have very strong intuition about the behavior of $p(t)$
based on our experience with objects moving in time. But we know that our intuition
can be wrong. For example suppose $p(t)$ gives the position of a
vertically thrown ball and that $\dfdxat{p}{t}{\tau}=10
\frac{\text{feet}}{\text{second}}$. The ball is thus moving upward
very quickly so we would expect that it will continue moving
upward. But if the ball is thrown indoors and $\tau$ is the instant
when the ball strikes the ceiling then of course it does not continue
rising. Counter to our intuition it bounces and immediately moves
downward.

Throughout\notetoself{This would be a good place to foreshadow
  continuity.} this discussion the word ``interval'' kept popping up
because it is more useful to know the behaviour of a function on an
interval than at a single point.  We would rather have a version of
Definition~\ref{def:first-deriv-test} that guarantees our function is
increasing \emph{on an interval} rather than at a single point. The
problem with that is that we already know what we mean by ``increasing
on an interval.'' It means that if $f(x)$ is increasing on some
interval $I$ this means that if $x, y\in I$ and $x\ge y$ then
$f(x)\ge f(y).$ So a test that does that will have to be a theorem
that gives conditions to be met and a proof that if those conditions
are met then $f(x)$ is increasing on the interval.

When we return to this topic in Section~\ref{sec:first-deriv-test} we
will prove exactly such a theorem.
\end{Digression}

\begin{ProblemSection}
  \begin{embeddedproblem}{}
    Show that at $x\approx-3.45$ the graph of the function
    $y(x)=\frac{x(x-1)}{x^3-27}$ has:
    \begin{enumerate}[label={\bf{}(\alph*)}]
    \item Both a global and a local minimum on the interval
      $(-5,0)$.
    \item Only a local minimum on the interval $\left(-5,
        \frac53\right)$. 
    \item Does the graph of this function have a global maximum on
      either interval? Explain.
    \end{enumerate}
  \end{embeddedproblem}

  \begin{embeddedproblem}{}
    Does $f(x)=\sin(x)$ have a global extremum (minimum or maximum) on
    the interval $\left(0,\frac{\pi}{2}\right)$? How about on $(0,
    \pi)$? On $(0, 2\pi)$? Explain.
  \end{embeddedproblem}


  \begin{embeddedproblem}{}
    Suppose that $O(x) = \frac{x^2}{1+x^2}$ and that the domain of the
    problem is all real numbers ($\RR$).
    \begin{enumerate}[label={\bf{}(\alph*)}]
    \item Find all local extrema.
    \item Find all global extrema. 
    \end{enumerate}
  \end{embeddedproblem}

  \begin{embeddedproblem}{}
    \label{problem:WitchTri}
    We draw the right triangle whose hypotenuse starts at the origin,
    $(0,0)$ and ends on the Witch of Agnesi.  The other legs are as
    shown in the figure
    below:\\
    \centerline{\includegraphics*[height=1.9in,width=3.5in]{../Figures/MaxTri}}
    \begin{enumerate}[label={\bf{}(\alph*)}]
    \item Of all possible such triangles in the interval $[0,\infty)$ which
      one has the largest  area?\\
    \item Of all possible such triangles in the interval $(-\infty,0]$ which
      one has the largest  area?\\
    \end{enumerate}

    
    The fact is that there is no minimum area. As $x$ grows larger the
    area of our triangle, $A(x) = \frac{x}{2(x^2+1)}$ grows smaller. In
    fact we can make the area as small as we like by taking $x$ large
    enough

    How large does $x$ need to be to guarantee that:
    \begin{enumerate}[label={\bf{}(\alph*)}]
      \begin{multicols}{2}
      \item $A(x)<\frac14$
      \item $A(x)<\frac{1}{5}$
      \item $A(x)<\frac{1}{10}$
      \item $A(x)<\frac{1}{100}$
      \item Now suppose $\eps$ is a very small number. How large does $x$
        need to be to guarantee that $A(x)<\eps$?
      \end{multicols}
    \end{enumerate}


  \end{embeddedproblem}

  \begin{embeddedproblem}{}
    These problems are slight variations on
    example~\ref{problem:WitchTri} in that the graph of the function
    has essentially the same shape.  Proceed as in the example above
    to find the area of the triangle with maximum area for each of the
    functions given.
    \begin{enumerate}[label={\bf{}(\alph*)}]{}
    \item $f(x) = \frac{k}{x^2+1}$, where $k$ is any {\bf{}k}onstant.
    \item $f(x) = \frac{1}{x^2+3}$
    \item $f(x) = \frac{1}{5x^2+1}$
    \item $f(x) = \frac{a}{bx^2+c}$, where $a$, $b$, and $c$ are
      arbitrary constants.
    \end{enumerate}
  \end{embeddedproblem}

  \begin{embeddedproblem}{}
    This problem is a more substantial variation on the same theme:
    Build a triangle in the same manner as in
    example~\ref{problem:WitchTri} but using the function
    $f(x)=\frac{x}{x^2+1}$ in place of the Witch. Find the coordinates
    of the triangle with maximum area, and find the maximum area.
  \end{embeddedproblem}

  \begin{embeddedproblem}
    Construct a triangle in the same way we did in
    problem~\ref{problem:WitchTri}, but in place of the Witch use the
    function $f(x)=\frac{1}{\sqrt{x}}$. Does the area of the triangle
    get smaller as $\alpha$ moves farther to the right\notetoself{We
      should refer back to this section when we get to L'H\^o{}pital's
      Rule-- ecb -- Aug 1, 2020}?
  \end{embeddedproblem}

  \begin{embeddedproblem}{}
    Suppose we construct a triangle in the
    same way as before, but this time we use the function
    $$
    f(x)=5x^4+3x^3-7x^2-5x+6
    $$
    instead of the Witch. In this case our sketch looks like this:\\
    \centerline{\includegraphics*[height=2in,width=3.5in]{../Figures/MaxTri3}}
    \begin{enumerate}[label={\bf{}(\alph*)}]
    \item Of all possible such triangles in the interval $[0,\infty)$ which
      one has the largest  area?\\
    \item Of all possible such triangles in the interval $(-\infty,0]$ which
      one has the largest  area?\\
    \end{enumerate}
  \end{embeddedproblem}

  \begin{embeddedproblem}{}
    Consider the polynomial
    $$
    y = (x+1)^3(x-1)(x+3).
    $$
    \begin{description}
    \item[(a)] Show that $\dfdxat{y}{x}{x=-1}=0.$ \hint{You can save
        yourself a lot of work on this problem by keeping you eye on the
        goal. You are not asked to compute $\dfdx{y}{x}.$ The problem is to
        show that $\left.\dfdx{y}{x}\right|_{x=-1}=0.$ You don't need to
        fully expand $\dfdx{y}{x}$ to do that.}
    \item[(b)]  Graph $y$ near $x=1$ to convince yourself that $y$ is
      neither a maximum nor a minimum at $x=-1$.
    \item[(c)] We don't really need a polynomial as complicated as the one
      given in part (a) to make the point we're getting at. Repeat parts
      (a) and (b) with the polynomial $y=x^3$ at $x=0.$
    \end{description}
  \end{embeddedproblem}
\end{ProblemSection}

\subsection{Optimization and Concavity}
\label{sec:second-deriv-test}
Calling definition~\ref{def:first-deriv-test} \emph{First} Derivative
Test suggests that there might be something called a \emph{Second}
Derivative Test, and indeed there is. Just as the First Derivative
Test tells us whether the graph of a function is \term{increasing} or
\term{decreasing}, the Second Derivative Test tells us whether the
graph is \term{concave upward} or \term{concave downward}.


\subsection*{\underline{Concavity}}
Clearly this will not make sense until we've understood what ``concave
upward'' and ``concave downward'' mean. The following examples are
illuminative.
Each of the following graphs is  ``concave upward,''\\
\centerline{\includegraphics*[height=1in,width=2in]{../Figures/ConcaveUp}}
whereas the following graphs are ``concave downward.''\\
\centerline{\includegraphics*[height=1in,width=2in]{../Figures/ConcaveDown}}
Obviously concave upward means ``like a bowl that opens up,'' and concave downward
means ``like a bowl that opens down. However it is easy to get the wrong
impression from simple examples so let's look at some more.
The graph of the natural exponential and the natural logarithm\\
\centerline{\includegraphics*[height=1.5in,width=1.5in]{../Figures/LogAndExp}}
are  concave upward and concave downward, respectively despite the facts that
the exponential never turns up on the right, and the logarithm never
turns down on the right.

Do you think this graph is concave upward or concave downward?\\
\centerline{\includegraphics*[height=1.2in,width=2in]{../Figures/ConcaveUpDown}}


Clearly it is both. On the interval $(0,\pi)$ it is concave downward and
on the interval $(\pi,2\pi)$ it is concave upward. Concavity is like
increasing/decreasing in that we will always need to specify the
interval on which a given function is either concave upward or concave downward.

Now what about these two?\\
\centerline{\includegraphics*[height=2in,width=4in]{../Figures/ConcaveUpDownUnclear}}

The graph on the right is clearly\footnote{Never mind that we can't
  actually see \emph{all} of the graph.}  concave upward when $x<0,$ but
what about when $x>0?$ Is it concave upward there? Concave downward?  Neither?
Both? It is very hard to tell from the graph.  The one on the left is
similarly vague in the region around $x=0.$

If it is important for us to distinguish concave upward from concave
downward\footnote{It is.}, clearly we will need a more precise method than
simply looking at graphs. We'll proceed slowly.  Since we're studying
Calculus it seems likely that the derivative will play an instrumental
role, so let's see what we can say about the derivative of a function
whose graph is concave upward.

Consider our first example $y(x) = x^2+1.$ This is concave upward
everywhere so we will consider how its derivative changes as $x$ 
increases. The following sketches show the graph and
its tangent at $x=0,$ $x=2,$ and $x=3.$\\
\centerline{\includegraphics*[height=3in,width=5.5in]{../Figures/ConcaveUpWithTangent}}
Notice that as $x$ increases the tangent line gets steeper and
steeper. That is to say its slope gets larger.  More formally,
the derivative increases.

\begin{ProficiencyDrill}
  Look back upward to the graph of the other ``concave up'' examples and
  convince yourself that this is true of all of them: As $x$ increases
  the derivative of each function also increases.
\end{ProficiencyDrill}

Does that mean that the derivative of a concave downward function will
decrease as $x$ increases? Yes, of course it does.

\begin{ProficiencyDrill}
  Look back up to the graph of the  ``concave down'' examples and
  convince yourself that this is true of all of them: As $x$ increases
  the derivative of each function decreases.
\end{ProficiencyDrill}

So, apparently in order to distinguish where a function is concave
upward from where it is concave downward we need to locate the
intervals where the derivative is increasing and where it is
decreasing.  So we need a tool to tell us when a function is
increasing and when it is decreasing.

But we already have that tool! It's called the  First Derivative Test!

To find out
where the graph of $y(x)$ is concave upward or concave downward, we ask where
its derivative, $\dfdx{y}{x}$, is increasing and where it is
decreasing. To find out where $\dfdx{y}{x}$ is increasing and
decreasing we ask where \emph{its} derivative, $\dfdxn{y}{x}{2}$, is
positive and negative

By the First Derivative Test\footnote{Because we are applying the
  First Derivative Test to the first derivative the language gets a
  little complex here. Read it carefully.} the derivative
$\dfdx{y}{x}$ will be increasing (so $y(x)$ is concave
upward) wherever \alert{its} derivative is positive
$\left(\dfdxn{y}{x}{2}>0\right)$ and will be decreasing (so $y(x)$ is
concave downward) wherever its derivative is negative
$\left(\dfdxn{y}{x}{2}<0\right)$.

So, to determine where the graph of $y(x)$ is concave upward or concave
down we need to look at the \emph{second derivative,}
$\dfdxn{y}{x}{2},$ and ask where it is positive and where it is
negative.

More formally:

\begin{mytheorem}[Concavity]
  \label{thm:what-second-deriv}
  Suppose $$y=f(x)$$ where $f$ is some function  of $x.$
  The following statements are true:
  \begin{enumerate}
  \item If $\eval{\dfdxn{y}{x}{2}}{x}{a}>0,$ then $\dfdx{y}{x}$ is
    increasing, and therefore $f$ is concave upward, at
    $x=a$. 
  \item If $\eval{\dfdxn{y}{x}{2}}{x}{a}<0,$ then $\dfdx{y}{x}$ is
    decreasing, and therefore $f$ is concave downward at
    $x=a$. 
  \item If $\eval{\dfdxn{y}{x}{2}}{x}{a}=0$ then $\dfdx{y}{x}$ is
    neither increasing nor decreasing, and therefore $f(x)$ is neither
    concave upward nor concave downward at $x=a$.
  \end{enumerate}
\end{mytheorem}

Do you see why Theorem~\ref{thm:what-second-deriv} is so similar to
Theorem~\ref{thm:FDT}? Theorem~\ref{thm:what-second-deriv} uses the
conclusions of Theorem~\ref{thm:FDT}, applied to the \emph{first
  derivative $f$} to draw logical conclusions about the concavity of
the graph of $f.$ So naturally the conclusion have the same form.

A couple of simple examples will illustrate this theorem.
\begin{myexample}
  Let $y=x^2.$ Then $\dfdx{y}{x}=2x,$ and $\dfdxn{y}{x}{2}=2.$ Since
  the second derivative of $y$ is \emph{always} positive the graph
  must always be concave upward.

  \begin{ProficiencyDrill}
    Is it? If you don't know graph it and see.
  \end{ProficiencyDrill}
\end{myexample}

\begin{myexample}
  Let $y=x^3.$ Then $\dfdx{y}{x}=3x^2,$ and $\dfdxn{y}{x}{2}=3x.$ Since
  the second derivative of $y$ is negative on the interval
  $(-\infty,0)$ the graph
  must be concave downward on $(-\infty, 0)$. Since the second derivative of $y$ is
  positive on the interval $(0,\infty)$ the graph must be concave upward
  on $(0,\infty)$.

  \begin{ProficiencyDrill}
    Is it? If you don't know, graph it and see.
  \end{ProficiencyDrill}
\end{myexample}

As we saw when we were discussing the First Derivative Test, 
the intervals where a function is increasing or decreasing  are
precisesly the intervals between the possible transition points.

Thus it follows that to find the intervals where $y(x)$ is
always concave upward or concave downward we will need to find the possible
transition points of the derivative; those points where the second
derivative of $y$, is zero or undefined. 

\begin{embeddedproblem}
  Suppose $y=x^4-x^2.$
  % Then $\dfdx{y}{x}=4x^3-2x$, and
  % $\dfdxn{y}{x}{2} =  12x^2-2$. To find the possible transition points
  % for \emph{concavity} we set   $\dfdxn{y}{x}{2} =  12x^2-2 =0$ and
  % solve. This yields, $x=\frac{1}{\sqrt{6}}$ and
  % $x=-\frac{1}{\sqrt{6}}$

  \begin{enumerate}[label={\bf{}(\alph*)}]
  \item Show that the graph of $y$ is concave upward  on the intervals
    $(-\infty,-1/\sqrt{6})$ and $(1/\sqrt{6},\infty)$.
  \item Show that the graph of $y$ is concave downward on  the interval
    $(-1/\sqrt{6},1/\sqrt{6})$.
  \end{enumerate}
  % \begin{enumerate}{}
  % \item Show that $\left.\dfdxn{y}{x}{2}\right|_{x=0}< 0.$
  % \item Sketch the graph of $f$ and confirm visually that it is
  %   concave downward at $x=0.$
  % \end{enumerate}
\end{embeddedproblem}


This last problem also displays a most useful fact about the second
derivative. Notice that at $ x=0$ the \emph{first} derivative of
$y=x^4-x^2$ is zero, so the tangent line is horizontal and the second
derivative is negative, so the graph is concave downward.  Taking
these two facts together we conclude that $y(x)$ must have a (local)
maximum at $x=0.$

Similarly, given any curve $y,$ if $\left.\dfdx{y}{x}\right|_{x=a}=0$
and $\left.\dfdxn{y}{x}{2}\right|_{x=a}> 0$
then $y$ must have a (local) minimum at $x=a.$

We will formalize this into a theorem called the
Second Derivative Test:

\begin{mytheorem}[The Second Derivative Test]
  Suppose the graph of $y$ is some curve and that at some point $x=a$
  $$\left.\dfdx{y}{x}\right|_{x=a}=0.$$
  then the following statements are true:
  \begin{enumerate}[label={\bf{}(\alph*)}]{}
  \item If $\left.\dfdxn{y}{x}{2}\right|_{x=a}>0$ then $y$ has a local
    minimum at $x=a.$
  \item If $\left.\dfdxn{y}{x}{2}\right|_{x=a}<0$ then $y$ has a local
    maximum at $x=a.$
  \item If $\left.\dfdxn{y}{x}{2}\right|_{x=a}=0$ then we get
    \emph{no} information about $y.$
  \end{enumerate}
\end{mytheorem}

\begin{ProficiencyDrill}
  Do you see why we get no information in part (c)? Analyze the
  concavity of the graphs of each of the
  following functions to see how part (c) can be true.
  \begin{enumerate}[label={\bf{}(\alph*)}]
    \begin{multicols}{3}
    \item $f(x)=x^4$
    \item $f(x)=-x^4$
    \item $f(x)= x^3$
    \end{multicols}
  \end{enumerate}
\end{ProficiencyDrill}

%%% Local Variables: 
%%% mode: latex
%%% outline-minor-mode: t
%%% TeX-master: "DiffCalc"
%%% eval:(load-file "../macros.el")
%%% End: 



